
\chapter{主存、总线、I/O 与 DMA}

前面已经解释了最小 CPU 如何在内部取指、译码、执行和提交状态。若停在这里，读者只能得到一台”会在寄存器之间搬数”的小机器。真正的整机还需要大容量主存、可寻址设备、外部事件、等待、错误和批量数据搬运。本章补上这条桥：CPU 怎样把一个地址变成一次存储器或设备事务，设备怎样把完成事件通知 CPU，DMA 又怎样在不让 CPU 逐字搬运的情况下读写主存。

\begin{keyidea}
整机硬件不是“CPU 加一块内存”的粗线图，而是由一组定义明确的事务组成。每次访问都必须说明地址、方向、字节选择、目标片选、握手、完成、错误和副作用。只有这些边界明确，CPU 才能可靠地取指、访存、访问设备和处理中断。
\end{keyidea}

\begin{pdfdiagram}
  \node[hw state, minimum width=1.55cm] (cpu) at (0,0) {CPU};
  \node[hw compute, minimum width=1.95cm] (dec) at (2.35,0) {地址译码};
  \node[hw state, minimum width=1.65cm] (rom) at (5.0,1.0) {ROM};
  \node[hw state, minimum width=1.65cm] (ram) at (5.0,0) {SRAM};
  \node[hw control, minimum width=1.65cm] (io) at (5.0,-1.0) {MMIO};
  \node[hw control, minimum width=1.75cm] (irq) at (7.45,-1.0) {中断/DMA};
  \draw[hw bus] (cpu) -- node[above,hw label]{地址/控制} (dec);
  \draw[hw bus] (dec.east) |- (rom.west) node[pos=0.25,above,hw label]{片选};
  \draw[hw bus] (dec) -- (ram);
  \draw[hw bus] (dec.east) |- (io.west);
  \draw[hw return] (io) -- (irq);
  \draw[hw return] (irq.south) -- ++(0,-0.5) -| (cpu.south) node[pos=0.35,below,hw label]{完成/事件};
  \node[hw label, text width=8.8cm] at (3.75,-2.45) {整机事务从 CPU 发起，但完成、等待和错误由目标与互连共同决定；设备还可通过中断或 DMA 把事件带回 CPU。};
\end{pdfdiagram}
\diagramnote{从最小 CPU 到整机事务。地址译码决定目标，握手和事件决定本次访问何时真正完成。}

\section{一、主存不是抽象数组}

软件把内存看成按地址编号的字节序列，硬件却必须用地址线、译码器、存储阵列、数据线和控制信号实现这件事。一次读访问至少包含四个事实：发起者给出地址，译码逻辑选中目标，目标经过内部访问后驱动数据，发起者在规定边界采样。一次写访问则要求地址、写数据、字节使能和写控制在目标采样窗口内稳定。

\begin{longtable}{L{0.18\linewidth}L{0.34\linewidth}L{0.38\linewidth}}
\toprule
对象 & 硬件含义 & 关键问题 \\
\midrule
地址 & 指出本次事务目标位置 & 高位译码是否只选中一个目标 \\
数据 & 读时由目标返回，写时由发起者提供 & 同一时刻是否只有一个驱动者 \\
读写控制 & 区分本次事务方向和类型 & 输出使能和写使能是否互斥 \\
字节使能 & 指明哪些 8-bit 车道有效 & 字节写不会破坏相邻数据 \\
ready/valid & 指明请求或响应是否真正完成 & 慢目标能否插入等待而不丢地址和数据 \\
\bottomrule
\end{longtable}

\topic{ROM、SRAM 和 DRAM 的读写不是同一种动作}
ROM 或 Flash 适合保存启动代码，掉电后内容仍在；SRAM 接口直观，地址稳定后在访问时间内返回数据，适合教学板、cache 和小容量片上 RAM；DRAM 密度高，但单元需要刷新，访问还要经过 bank、行、列、预充电和控制器调度。把这些对象都叫“内存”没有错，但在硬件设计中必须区分它们的时序和控制边界。

\begin{longtable}{L{0.18\linewidth}L{0.34\linewidth}L{0.38\linewidth}}
\toprule
存储类型 & 典型用途 & 关键边界 \\
\midrule
ROM/Flash & 复位入口、固件、常量表 & 复位地址必须映射到有效内容，读时序满足取指 \\
SRAM & cache、实验板 RAM、小容量片上存储 & \code{CS}/\code{OE}/\code{WE} 互斥，地址和数据满足建立时间 \\
DRAM & 大容量主存 & 控制器负责 activate、read/write、precharge、refresh \\
寄存器文件 & CPU 内部近距离状态 & 多读端口、写端口、旁路和同周期读写规则明确 \\
\bottomrule
\end{longtable}

\topic{异步 SRAM 是最好的第一块主存}
教学整机可以先使用低速异步 SRAM。读周期中，地址和片选稳定后，打开输出使能，SRAM 在访问时间后驱动数据线；写周期中，输出使能关闭，发起者驱动数据线，在写使能有效窗口内把数据写入目标字节。这个模型足够简单，也足以暴露硬件最常见错误：片选重叠、输出争用、写使能过早、字节使能错误。

\begin{lstlisting}
SRAM read:
  address stable
  CS active, OE active, WE inactive
  wait access time
  sample data

SRAM write:
  address and write_data stable
  CS active, OE inactive, WE active for write window
  selected byte lanes store data
\end{lstlisting}

\begin{pdfdiagram}[x=1cm,y=1cm]
  % ADDR — 总线由无效翻转到有效
  \draw[BookTeal,line width=0.85pt] (0.6,1.7) -- (1.6,1.2) -- (8.4,1.2);
  \draw[BookTeal,line width=0.85pt] (0.6,1.2) -- (1.6,1.7) -- (8.4,1.7);
  \node[hw label, anchor=east] at (0.5,1.45) {ADDR};
  \node[hw label] at (5.0,1.45) {地址稳定};
  % CS_n/OE_n — 低有效: 高→低→高
  \draw[BookAccent,line width=0.85pt] (0.6,0.55) -- (1.6,0.55) -- (1.6,0.05) -- (7.6,0.05) -- (7.6,0.55) -- (8.4,0.55);
  \node[hw label, anchor=east] at (0.5,0.3) {$\overline{CS}$/$\overline{OE}$};
  % DATA — 高阻(中间虚线) → 有效数据 → 释放回高阻
  \draw[BookMuted,line width=0.65pt,dashed] (0.6,-0.75) -- (4.6,-0.75);
  \draw[BookMuted,line width=0.65pt,dashed] (7.6,-0.75) -- (8.4,-0.75);
  \draw[BookTeal,line width=0.85pt] (4.6,-0.45) -- (7.6,-0.45) -- (7.6,-1.05) -- (4.6,-1.05) -- cycle;
  \node[hw label, anchor=east] at (0.5,-0.75) {DATA};
  \node[hw label] at (2.6,-0.4) {高阻};
  \node[hw label] at (6.1,-0.75) {数据有效};
  % 区间竖线
  \draw[BookRule,line width=0.45pt,dashed] (1.6,2.15) -- (1.6,-1.3);
  \draw[BookRule,line width=0.45pt,dashed] (4.6,2.15) -- (4.6,-1.3);
  \draw[BookRule,line width=0.45pt,dashed] (7.6,2.15) -- (7.6,-1.3);
  \draw[BookMuted,line width=0.6pt,<->] (1.6,2.0) -- node[hw label, above]{访问时间（等待）} (4.6,2.0);
  \draw[BookMuted,line width=0.6pt,<->] (4.6,2.0) -- node[hw label, above]{CPU 采样窗口} (7.6,2.0);
  \node[hw label, text width=8.6cm] at (4.5,-1.7) {等待状态的目的，是把采样点推迟到数据有效之后，同时保持地址、片选和方向不变。};
\end{pdfdiagram}
\diagramnote{异步 SRAM 读周期：地址先稳定，$\overline{CS}$/$\overline{OE}$ 拉低后存储器需要一段访问时间才把数据驱动到总线上；CPU 只能在数据有效后采样。三条线各有明确的翻转时刻——读时序图读的就是这些边沿，而不是三条平直线。}

\topic{SRAM 写周期的驱动权与写窗口}
写周期与读周期使用同一组地址线、数据线和片选，差别在数据线的驱动方向和写使能窗口。读周期中 $\overline{OE}$ 有效，由 SRAM 驱动数据线，发起者只负责在数据有效后采样；写周期中 $\overline{OE}$ 必须保持无效，SRAM 的输出驱动器关闭，改由发起者把写数据驱动到数据线上——同一组数据线在任何时刻只能有一个驱动者。写使能 $\overline{WE}$ 拉低形成写窗口，规格书围绕它规定三条边界：地址必须在窗口打开前已经稳定，写数据必须在窗口关闭前满足建立时间，窗口关闭后还要继续保持一段保持时间，窗口本身也有最小宽度。这些要求来自同一个事实：存储阵列在 $\overline{WE}$ 回升、窗口关闭的边沿附近把数据线上的值写入选中单元，窗口内任何一次地址或数据翻动，都可能把错误的值写进错误的单元。因此写周期的固定顺序是：先给地址和数据，再打开写窗口，先关窗口，最后才撤数据。

\begin{waveform}[x=0.85cm,y=0.95cm]
  % 区间竖线：x=2 地址稳定, x=4 窗口打开, x=9 窗口关闭(写入边沿), x=9.8 数据释放
  \draw[BookRule, line width=0.45pt, dashed] (2,8.35) -- (2,0.25);
  \draw[BookRule, line width=0.45pt, dashed] (4,8.35) -- (4,0.25);
  \draw[BookRule, line width=0.45pt, dashed] (9,8.35) -- (9,0.25);
  \draw[BookRule, line width=0.45pt, dashed] (9.8,8.35) -- (9.8,0.25);
  % ADDR — 窗口打开前先稳定，窗口关闭后才允许变化
  \draw[BookTeal, line width=0.9pt] (0.5,6.5) rectangle (2,7.2);
  \node at (1.25,6.85) {旧值};
  \draw[BookTeal, line width=0.9pt] (2,6.5) rectangle (12,7.2);
  \node at (7,6.85) {地址稳定};
  \draw[BookMuted, line width=0.65pt, dashed] (12,6.85) -- (13,6.85);
  \node[hw label, anchor=east] at (0.3,6.85) {ADDR};
  % CS_n — 低有效: 地址稳定后拉低
  \draw[BookAccent, line width=0.9pt] (0.5,5.6) -- (2.5,5.6) -- (2.5,4.9) -- (10.5,4.9) -- (10.5,5.6) -- (13,5.6);
  \node[hw label, anchor=east] at (0.3,5.25) {$\overline{CS}$};
  % WE_n — 写窗口
  \draw[BookAccent, line width=0.9pt] (0.5,4.2) -- (4,4.2) -- (4,3.5) -- (9,3.5) -- (9,4.2) -- (13,4.2);
  \node[hw label, anchor=east] at (0.3,3.85) {$\overline{WE}$};
  % OE_n — 全程保持无效
  \draw[BookAccent, line width=0.9pt] (0.5,2.8) -- (13,2.8);
  \node[hw label, anchor=east] at (0.3,2.8) {$\overline{OE}$};
  \node[hw label] at (6.75,2.42) {全程保持无效};
  % DATA — 高阻(中间虚线) → 发起者驱动写数据 → 释放回高阻
  \draw[BookMuted, line width=0.65pt, dashed] (0.5,1.4) -- (4.5,1.4);
  \draw[BookMuted, line width=0.65pt, dashed] (9.8,1.4) -- (13,1.4);
  \draw[BookTeal, line width=0.9pt] (4.5,1.0) -- (9.8,1.0) -- (9.8,1.8) -- (4.5,1.8) -- cycle;
  \node[hw label, anchor=east] at (0.3,1.4) {DATA};
  \node[hw label] at (2.5,1.78) {高阻};
  \node at (7.15,1.4) {发起者驱动：写数据稳定};
  \node[hw label] at (11.4,1.78) {高阻};
  % 顶部标注: 地址先行 / 写窗口
  \draw[BookMuted, line width=0.6pt, <->] (2,7.95) -- node[hw label, above]{地址先稳定} (4,7.95);
  \draw[BookMuted, line width=0.6pt, <->] (4,7.95) -- node[hw label, above]{写窗口} (9,7.95);
  % 底部标注: 建立 / 保持
  \draw[BookMuted, line width=0.6pt, <->] (4.5,0.55) -- node[hw label, below]{数据建立} (9,0.55);
  \draw[BookMuted, line width=0.6pt, <->] (9,0.55) -- node[hw label, below]{数据保持} (9.8,0.55);
\end{waveform}
\diagramnote{异步 SRAM 写周期：地址与片选先稳定，$\overline{OE}$ 全程无效，发起者在写窗口内驱动数据线。存储阵列在 $\overline{WE}$ 回升、窗口关闭的边沿把数据线上的值写入选中单元，因此数据必须在窗口关闭前满足建立时间、关闭后满足保持时间，地址在整个窗口内不能翻动——与读周期对照，差别就在数据线由谁驱动、哪个边沿是真正的写入点。}

\topic{DRAM 需要控制器而不是直接接地址线}
DRAM 的容量来自更密集的存储单元，代价是接口复杂。控制器要先打开某个 bank 的某一行，再读写列；换行前需要预充电；即使没有普通访问，也要周期性刷新；高速 DRAM 还需要训练采样窗口。CPU 看到的是“主存读写”，但实际路径经过控制器、队列、bank 调度、行命中或行冲突。这也是为什么现代 CPU 必须靠 cache 把常用数据放近。

\begin{longtable}{L{0.2\linewidth}L{0.34\linewidth}L{0.36\linewidth}}
\toprule
DRAM 动作 & 硬件含义 & 时延来源 \\
\midrule
Activate & 打开某个 bank 中的一行 & 行译码和感放建立 \\
Read/Write & 在已打开行内选择列传输数据 & 列访问、突发传输和总线排队 \\
Precharge & 关闭当前行，准备访问另一行 & 位线恢复到可再次访问状态 \\
Refresh & 周期性恢复存储单元电荷 & 控制器必须暂停或调度普通访问 \\
\bottomrule
\end{longtable}

\topic{存储芯片的位宽扩展与容量扩展}
单片存储芯片的位宽和容量都有限，整机用两组方向相反的手段补齐。位宽扩展解决“字不够宽”：把若干片芯片并接同一组地址线、片选和读写控制，每片只接数据线的一个字节车道——四片 x8 SRAM 并接后，一次 32-bit 访问四片同时动作，各自收发 D[7:0]、D[15:8]、D[23:16]、D[31:24]，对外表现为一片 32-bit 宽的 SRAM。容量扩展解决“行不够多”：数据线和低位地址全部并接，新增的高位地址经译码器产生多路互斥片选，任一地址最多选中一片，行数随片数倍增。两种扩展经常同时使用，但边界必须分清：位宽扩展中各片同时被选中、各管一段数据线；容量扩展中各片互斥选中、数据线全部共用。混淆二者的直接后果，是两片同时驱动同一个字节车道的总线争用。

\begin{pdfdiagram}
  \node[hw control, minimum width=2.1cm] (src) at (0.7,0) {同一组地址\\$\overline{CS}$/$\overline{WE}$};
  \node[hw memory, minimum width=2.0cm] (c0) at (4.15,2.7) {x8 SRAM 片 0};
  \node[hw memory, minimum width=2.0cm] (c1) at (4.15,0.9) {x8 SRAM 片 1};
  \node[hw memory, minimum width=2.0cm] (c2) at (4.15,-0.9) {x8 SRAM 片 2};
  \node[hw memory, minimum width=2.0cm] (c3) at (4.15,-2.7) {x8 SRAM 片 3};
  \node[hw state, minimum width=2.2cm, minimum height=6.6cm] (dbus) at (7.9,0) {32-bit 数据\\D[31:0]};
  % 地址/控制共享干线 + 分层支路
  \draw[hw bus, -] (src.east) -- (2.5,0);
  \draw[hw bus, -] (2.5,2.7) -- (2.5,-2.7);
  \draw[hw bus, shorten <=0pt] (2.5,2.7) -- (c0.west);
  \draw[hw bus, shorten <=0pt] (2.5,0.9) -- (c1.west);
  \draw[hw bus, shorten <=0pt] (2.5,-0.9) -- (c2.west);
  \draw[hw bus, shorten <=0pt] (2.5,-2.7) -- (c3.west);
  % 各片各占一个字节车道
  \draw[hw bus] (c0.east) -- node[above,hw label]{D[7:0]} (6.8,2.7);
  \draw[hw bus] (c1.east) -- node[above,hw label]{D[15:8]} (6.8,0.9);
  \draw[hw bus] (c2.east) -- node[above,hw label]{D[23:16]} (6.8,-0.9);
  \draw[hw bus] (c3.east) -- node[above,hw label]{D[31:24]} (6.8,-2.7);
  \node[hw label, text width=2.2cm] at (10.3,0) {容量扩展则靠片选译码增加行数};
\end{pdfdiagram}
\diagramnote{位宽扩展：四片 x8 SRAM 并接同一组地址、$\overline{CS}$ 和 $\overline{WE}$，同时动作、各管一个字节车道，对外是一片 32-bit SRAM。容量扩展方向相反：数据线全部并接，高位地址译码出互斥片选来增加行数，任一地址最多一片被选中。}

\topic{DRAM 单元用电容存电荷，读出后必须重写}DRAM 的一个 bit 由一个访问晶体管和一个存储电容构成（1T1C）。位线先预充到 $V_{CC}/2$；字线打开晶体管后，电容与位线共享电荷，位线电压略升或略降，读出放大器（感放）判别方向并输出 0/1。两个后果决定了 DRAM 的全部特殊性：其一，读出是破坏性的，感放必须把满摆幅电平经位线写回单元，否则原值丢失；其二，电容持续漏电，电荷在几十毫秒内就衰减到无法判别，控制器必须周期性刷新——逐行打开并重写。规格书里的 Activate、Precharge、Refresh 和刷新周期（典型要求 64 ms 内把全部行刷新一遍），都直接来自这两件事。

\begin{pdfdiagram}
  % 位线
  \draw[hw bus] (0.4,2.2) -- (4.9,2.2);
  \node[hw label, above] at (0.4,2.2) {位线 BL（先预充到 $V_{CC}/2$）};
  % 访问晶体管
  \pic at (1.6,1.55) {hw nmos};
  \draw[hw bus] (1.6,1.97) -- (1.6,2.2);
  \draw[hw bus] (0.4,1.55) -- (0.95,1.55);
  \node[hw label, above] at (0.4,1.55) {字线 WL};
  % 存储电容
  \draw[hw bus] (1.6,1.13) -- (1.6,0.92);
  \draw[hw bus] (1.38,0.92) -- (1.82,0.92);
  \draw[hw bus] (1.38,0.78) -- (1.82,0.78);
  \draw[hw bus] (1.6,0.78) -- (1.6,0.5);
  \pic at (1.6,0.5) {hw gnd};
  \node[hw label, anchor=west] at (2.0,0.85) {存储电容 C};
  % 读出放大器
  \node[hw compute, minimum width=3.0cm] (sa) at (6.5,2.2) {读出放大器（感放）\\判别后满摆幅回写};
  \node[hw label, anchor=west] at (8.15,2.2) {读出的 0/1};
  \node[hw label, text width=9.6cm] at (4.3,-0.75) {打开字线后，电容与位线共享电荷，位线电压轻微上升或下降；感放判别方向输出 0/1，并立刻把满摆幅电平写回单元——读出即重写。};
\end{pdfdiagram}
\diagramnote{1T1C 存储单元与读出放大器。刷新并不是额外功能，而是控制器替所有行周期性重复这个“读出—回写”过程；这也是 DRAM 不能像 SRAM 那样直接把地址线接上去用的原因。}

\topic{DRAM 的行列地址复用与突发传输}
DRAM 的地址引脚不随容量等比增长，靠的是行列地址复用：同一组地址引脚先送行地址、再送列地址。一个 $2^{20}$ 字的阵列直接接地址需要 20 根地址线，复用后 10 根就够——$\overline{RAS}$ 有效时把引脚上的值锁存为行地址，$\overline{CAS}$ 有效时把同一组引脚上的值锁存为列地址。省下的引脚直接降低封装成本，代价是接口从“给地址等数据”变成一串带时序要求的命令。行地址译码后打开整行，整行数据被感放读出并留在行缓冲中；此后的列访问只是从行缓冲里选一段送到数据总线。突发传输利用的正是这一点：给一个起始列地址，DRAM 按内部列计数器连续多拍输出相邻数据，后续各拍不再付出行打开的代价。行缓冲命中是突发便宜的根本原因，也是内存控制器尽量把相邻访问调度到同一行的理由。

\begin{pdfdiagram}
  \node[hw state, minimum width=2.0cm] (addr) at (0.6,0) {地址总线};
  \node[hw compute, minimum width=2.0cm] (mux) at (2.9,0) {行列复用器};
  \node[hw state, minimum width=1.9cm] (rowlatch) at (5.4,0.9) {行地址锁存};
  \node[hw state, minimum width=1.9cm] (collatch) at (5.4,-0.9) {列地址锁存};
  \node[hw compute, minimum width=1.4cm] (rowdec) at (7.7,0.9) {行译码};
  \node[hw compute, minimum width=1.4cm] (coldec) at (7.7,-0.9) {列译码};
  \node[hw memory, minimum width=2.0cm] (bank) at (10.2,0.9) {bank 阵列};
  \node[hw memory, minimum width=2.0cm] (rowbuf) at (10.2,-0.9) {行缓冲};
  \node[hw compute, minimum width=2.0cm] (burst) at (10.2,-2.4) {突发输出\\连续多拍};
  % 地址总线经共享干线分送到行、列锁存
  \draw[hw bus] (addr) -- (mux);
  \draw[hw bus] (mux.east) -- (5.4,0) -- (rowlatch.south);
  \draw[hw bus, shorten <=0pt] (5.4,0) -- (collatch.north);
  \draw[hw return] (5.4,1.55) -- (rowlatch.north);
  \node[hw label] at (5.4,1.85) {$\overline{RAS}$};
  \draw[hw return] (5.4,-1.55) -- (collatch.south);
  \node[hw label] at (5.4,-1.85) {$\overline{CAS}$};
  % 行路径与列路径
  \draw[hw bus] (rowlatch) -- node[above,hw label]{行} (rowdec);
  \draw[hw bus] (collatch) -- node[above,hw label]{列} (coldec);
  \draw[hw bus] (rowdec) -- (bank);
  \draw[hw bus] (coldec) -- (rowbuf);
  \draw[hw bus] (bank.south) -- node[right,hw label]{整行读入} (rowbuf.north);
  \draw[hw bus] (rowbuf.south) -- node[right,hw label]{按列连读} (burst.north);
  \node[hw label, text width=10.2cm] at (5.35,-3.5) {同一组地址引脚在 $\overline{RAS}$、$\overline{CAS}$ 两拍分别送出行、列地址；整行先读入行缓冲，列访问只是在缓冲里选段，行命中后的连续突发因此不必重开行。};
\end{pdfdiagram}
\diagramnote{行列地址复用与突发传输。$\overline{RAS}$ 拍锁存行地址并打开整行，$\overline{CAS}$ 拍锁存列地址；行缓冲命中后，连续列地址直接从行缓冲多拍输出，比重开一行便宜得多。}

行打开、列传输和预充电之间的最小间隔由一组时序参数规定，规格书把它们写成以时钟拍为单位的数字。DDR3-1600 的内部时钟为 800 MHz，一拍 1.25 ns；1600 MT/s 的数据速率来自同一时钟的双沿传输，时序参数仍按时钟拍计。常见标注 11-11-11-28 就是下表四个数字。

\begin{longtable}{L{0.12\linewidth}L{0.44\linewidth}L{0.16\linewidth}L{0.18\linewidth}}
\toprule
参数 & 含义 & DDR3-1600 & 纳秒换算 \\
\midrule
$t_{RCD}$ & Activate 打开行之后，到第一个列读/写命令的最小间隔 & 11 拍 & 13.75 ns \\
$t_{CL}$ & 列读命令发出后，到第一个数据出现在总线上（CAS 延迟） & 11 拍 & 13.75 ns \\
$t_{RP}$ & Precharge 关闭当前行之后，到允许再次 Activate 的最小间隔 & 11 拍 & 13.75 ns \\
$t_{RAS}$ & 一次 Activate 到允许 Precharge 的最小行打开时间 & 28 拍 & 35 ns \\
\bottomrule
\end{longtable}

换算很直接：拍数乘 1.25 ns。同一行内连续读，列延迟 13.75 ns 后就能紧跟突发数据；访问另一行则要先等 $t_{RP}$ 再重新 Activate，一次行冲突至少多花 $t_{RP}+t_{RCD}=22$ 拍，约 27.5 ns。这个数字就是内存调度器在乎行命中、软件在乎访问局部性的原因。

\topic{SDRAM 的突发长度与多 bank 交错}
Activate 打开一行之后，控制器可以对同一行连发多个列读写命令；每个列命令传的也不是一个数据，而是按约定的突发长度（burst length）连续传一串——DDR3 固定为 8（BL8），也允许切成 4。DDR 的“双倍数据率”指时钟的上升沿和下降沿各传一次数据，而不是把时钟翻倍：DDR3-1600 的内部时钟是 800 MHz，数据速率 1600 MT/s，所以一次 BL8 突发占 4 个时钟拍、交 8 个数据，这期间数据总线被这一串独占。

bank 是 DRAM 内部可独立 Activate、独立 Precharge 的子阵列，各有自己的行缓冲。一个 bank 在等 $t_{RCD}$、等数据的时候，命令总线可以对另一个 bank 发命令，把后者的行打开时间藏进前者的传输里，这就是多 bank 交错。沿用上一张表的 DDR3-1600 参数（1 拍 1.25 ns，$t_{RCD}=t_{CL}=t_{RP}=11$ 拍，$t_{RAS}=28$ 拍），比较“读两个不同行”在单 bank 与双 bank 下的时间轴：

\begin{lstlisting}
单 bank，两次访问落在同一 bank 的不同行:
  拍  0  ACT bank0 row0
  拍 11  RD  bank0 col0    (t_RCD 到期)
  拍 22  数据出现在总线，BL8 占用拍 22-25
  拍 28  PRE bank0         (t_RAS 到期才允许关行)
  拍 39  ACT bank0 row1    (t_RP 到期)
  拍 50  RD  bank0 col1    (t_RCD 到期)
  拍 61  第二组数据出现，共等 61 拍

双 bank 交错，两行分别在 bank0、bank1:
  拍  0  ACT bank0
  拍  1  ACT bank1         (不同 bank，行打开并行进行)
  拍 11  RD  bank0         (数据占拍 22-25)
  拍 15  RD  bank1         (数据占拍 26-29，恰好接上)
  拍 26  第二组数据出现，共等 26 拍
\end{lstlisting}

第二组数据从拍 61 提前到拍 26，省 35 拍、约 43.75 ns。省下的不是任何一条时序参数——单 bank 必须串行等完 $t_{RP}+t_{RCD}$ 才能再读，双 bank 让 bank1 的 Activate 与 bank0 的全过程并行，这段等待被完全覆盖；拍 15 才发第二条读命令，是因为数据总线仍由两个 bank 共享，bank0 的突发占着拍 22--25，bank1 的数据最早只能接在拍 26。这就是内存条标注 bank 数、控制器把相邻地址散到不同 bank、调度器重排请求的原因：行命中决定一次访问多快，bank 并行决定等待能不能被藏起来。这个例子是教学模型，为突出 bank 并行的收益，忽略了 $t_{RRD}$、$t_{FAW}$ 等对连续 Activate 的命令间约束。

\begin{longtable}{L{0.22\linewidth}L{0.32\linewidth}L{0.32\linewidth}}
\toprule
对照点 & 单 bank 顺序访问 & 双 bank 交错 \\
\midrule
第二组数据出现 & 拍 61 & 拍 26（省 35 拍，约 43.75 ns） \\
数据总线占用 & 拍 22-25、61-64，中间大量气泡 & 拍 22-29 连续无气泡 \\
关键等待 & 换行时 $t_{RP}+t_{RCD}$ 串行暴露 & 被另一 bank 的并行 Activate 覆盖 \\
检查点 & PRE 与 ACT 都卡在同一个 bank 上 & 命令错开，数据在共享总线上排队衔接 \\
\bottomrule
\end{longtable}

\topic{组相联 cache 的替换策略是精度与代价的取舍}
前面说过，主存的时延是结构性的，整机的应对是把热数据留在近处的 cache 里。本章后半部分会拆开 cache 的地址字段、写策略和一致性，这里先回答放进 cache 之后必然面对的问题：新行取回时组内已无空路，牺牲哪一路？直接映射没有这个问题，每个地址只有唯一槽位；组相联把它交给了替换策略。

精确 LRU（最久未使用）记录组内各路的年龄顺序，牺牲最老的一路，命中率最好，代价随路数急剧膨胀。2 路组相联时每组只需 1 位：这一位指出哪一路更久未用，每次命中取反即可。4 路要区分完整的年龄顺序，共 $4!=24$ 种排列，每组需 $\lceil\log_2 24\rceil=5$ 位，且每次命中都要重排年龄；8 路则需 $\lceil\log_2 40320\rceil=16$ 位。伪 LRU 把 4 路组织成两层二叉树，3 个内部节点各 1 位，每组 3 位，命中时沿路径翻转节点位；它选出的不保证是最老路，但命中率损失通常很小。FIFO 只记进入顺序，命中时不更新任何状态；随机替换连状态位都不要。

高组相联之所以普遍放弃精确 LRU，原因有两头：状态位按 $n!$ 增长还是小事，更难办的是每次命中都要更新年龄，路数越多更新逻辑的组合路径越长，直接拖慢命中时间；而 8 路、16 路的组里，“精确选最老”换来的命中率收益早已递减，随机替换与 LRU 的差距常常小到测量噪声量级。工程上的结论因此很清楚：2 路用 1 位精确 LRU，4 路用 3 位伪 LRU，更高组相联用伪 LRU 或随机。

\begin{longtable}{L{0.15\linewidth}L{0.25\linewidth}L{0.20\linewidth}L{0.28\linewidth}}
\toprule
策略 & 每组状态位（4 路） & 命中时是否更新 & 关键问题 \\
\midrule
精确 LRU & 5 位（$4!=24$ 种年龄序） & 每次命中重排年龄 & 路数增多后状态位和更新路径急剧膨胀 \\
伪 LRU & 3 位（两层树节点） & 沿树路径翻转节点位 & 选出的不保证是最老路，命中率略降 \\
FIFO & 2 位（轮流指针） & 不更新 & 不看访问热度，热行可能被轮到 \\
随机 & 0 位 & 不更新 & 实现最简，大组相联下命中率损失很小 \\
\bottomrule
\end{longtable}

\topic{硬件预取与软件预取都是拿带宽换命中}
替换策略决定牺牲谁，预取则试图让牺牲更少发生：顺序扫描数组时，用完一行 cache 行，下一行几乎马上要用，空间局部性让“将来要访问什么”变得可预测。硬件预取器监视 miss 流的地址规律：next-line 预取在用到第 $i$ 行时把第 $i+1$ 行取回；步长预取识别出固定步长后沿步长提前取回；预取流通常在 4 KiB 页边界停下——跨页地址可能未映射，为一次猜测触发缺页代价太大。

软件预取由指令显式给出提示。x86-64 提供 \code{prefetcht0}、\code{prefetchnta} 等预取指令：它们不改变任何架构可见状态，地址无效时也不触发缺页异常，只是提示硬件“这个地址很快会被读”。编译器或手写汇编常在循环里提前若干次迭代发出提示：

\begin{lstlisting}
; rdi 指向当前元素, rsi 是数组结尾, 每元素 8 字节
scan_loop:
    prefetcht0 [rdi + 512]   ; 提前 8 行(64 B 行)提示
    add     rax, [rdi]
    add     rdi, 8
    cmp     rdi, rsi
    jb      scan_loop
\end{lstlisting}

预取不是免费的：取错行会占住 cache 槽位、挤掉有用数据，每次预取本身也消耗内存带宽，过激的预取反而把真正需要的访问排在后面。可以粗算一笔账：顺序扫描大数组，行 64 B、元素 8 B，一行装 8 个元素；miss 时延按 200 拍、每个元素处理 4 拍计。无预取时每行一次 miss，每行花 $200+32=232$ 拍；预取把“取下一行”与“处理当前行”重叠，若提前量达到 $200\div 32\approx 7$ 行，处理完当前行时下一行已经在 cache 里，稳态每行只剩 32 拍，miss 被完全隐藏，吞吐提高约 $232/32\approx 7.3$ 倍。若预取只能提前一行，32 拍的处理盖不住 200 拍的取回，每行仍暴露约 168 拍。提前量必须匹配“处理一行的耗时”与“取回一行的耗时”之比，这正是预取器设计的核心数字。

\begin{longtable}{L{0.18\linewidth}L{0.34\linewidth}L{0.34\linewidth}}
\toprule
预取方式 & 触发时机 & 关键问题 \\
\midrule
next-line 硬件预取 & 用到第 $i$ 行时取第 $i+1$ 行 & 顺序流效果好，随机访问纯属浪费 \\
步长硬件预取 & 识别出固定地址步长后沿步长取 & 步长误判会持续污染 cache \\
软件预取指令 & 程序在用到数据之前显式提示 & 提示距离要折算成拍数，地址须确实会用 \\
跨页停止 & 预取流在 4 KiB 页边界停下 & 猜测性访问不得触发缺页副作用 \\
\bottomrule
\end{longtable}

\topic{上电自检的走步与棋盘格各查一类故障}
主存焊上板不等于可信：某位固定读成 0 或 1（stuck-at）、相邻数据线或地址线短路、写入后状态维持不住（数据保持故障），都可能让“写一个值再读回来”这种最简单的测试侥幸通过——它只验证了大部分位在此刻、对这组值是可写的。上电自检用图案化的写入-读回序列把故障类别分开：走步（walking 1/0）每次只让一位与众不同，棋盘格（checkerboard）让相邻位恒取反值。

走步 1 先写全 0 背景，再对同一单元依次写 \code{0x01}、\code{0x02}、……、\code{0x80}，每步读回比对，并复查其余单元仍是背景值。某位固定 0 会在轮到它走 1 时露馅，固定 1 会在背景检查中露馅；两位短路时走步值会变形——写 \code{0x01} 可能读回 \code{0x03}，直接指认相邻两位被短接。走步 0 把背景换成全 1、逐位走 0，对称覆盖另一半故障方向。棋盘格把 \code{0x55}（0101 0101）写满读回，再换成 \code{0xAA}（1010 1010）：任何相邻两位恒取反，相邻短路或位间耦合会让图案当场变形；保持故障则靠“写入后停一段再读”暴露。这些图案便宜、确定、能在无操作系统的裸机上运行，查不出所有故障，但恰好覆盖焊接与走线最常出的几类，因此是板级 bring-up 的固定节目。

以 4 字节小内存（addr0--addr3，8-bit 字节）为例，一次走步 1 加棋盘格的完整步骤如下。

\begin{longtable}{L{0.08\linewidth}L{0.38\linewidth}L{0.18\linewidth}L{0.24\linewidth}}
\toprule
步骤 & 操作 & 预期结果 & 检查点 \\
\midrule
1 & addr0--addr3 全写 \code{0x00}，逐字节读回 & 全部 \code{0x00} & 无固定 1，背景可写可读 \\
2 & addr0 依次写 \code{0x01}、\code{0x02}、……、\code{0x80}，每步读回 & 读回恒等于写入 & 每个数据位都能独立置 1 \\
3 & 步骤 2 每步之后读 addr1--addr3 & 仍为 \code{0x00} & 无地址串扰、无相邻线短路 \\
4 & addr1--addr3 重复步骤 2、3 & 同上 & 逐字节覆盖全部单元 \\
5 & 全写 \code{0x55} 读回，再全写 \code{0xAA} 读回 & 图案不变形 & 相邻位恒反，查相邻短路 \\
6 & 背景改 \code{0xFF}，逐位走 0（\code{0xFE}、\code{0xFD}、……） & 读回恒等于写入 & 对称覆盖固定 0 故障 \\
\bottomrule
\end{longtable}

\topic{校验位与 ECC 把位错从灾难降级为事件}
主存的位会错：宇宙射线、电压毛刺、保持时间失守都可能翻转一个 bit。最便宜的对策是奇偶校验——每字节加 1 个校验位，使连同校验位在内的 1 的个数恒为偶数（偶校验）。数据 \code{1011 0001} 已有 4 个 1，校验位记 0；读回时若某位翻转，1 的个数变奇，硬件立刻知道出错了。奇偶校验的边界同样明确：只能发现奇数个位错，两位同时翻转就漏检；即使发现，也不知道错在哪一位，只能报错，不能纠正。

汉明码用多个校验位换取定位能力。(7,4) 汉明码把 4 个数据位和 3 个校验位排成 7 位，位置从 1 编号；校验位放在编号为 2 的幂的位置（1、2、4），数据位占据 3、5、6、7。把位置编号写成 3 位二进制，第 $i$ 个校验位负责所有编号第 $i$ 位为 1 的位置。每个校验组各做一次偶校验，哪几组失败，其编号拼起来恰好就是出错位置的编号——这个把“哪些组错”翻译成“哪里错”的数称为校正子（syndrome）。

\begin{longtable}{L{0.12\linewidth}L{0.16\linewidth}L{0.28\linewidth}L{0.28\linewidth}}
\toprule
位置编号 & 存放内容 & 参与校验的校验位 & 说明 \\
\midrule
1 & $p_1$ & $p_1$ 自身 & 编号 001，最低位为 1 \\
2 & $p_2$ & $p_2$ 自身 & 编号 010 \\
3 & $d_1$ & $p_1$、$p_2$ & 编号 011 \\
4 & $p_4$ & $p_4$ 自身 & 编号 100 \\
5 & $d_2$ & $p_1$、$p_4$ & 编号 101 \\
6 & $d_3$ & $p_2$、$p_4$ & 编号 110 \\
7 & $d_4$ & $p_1$、$p_2$、$p_4$ & 编号 111 \\
\bottomrule
\end{longtable}

算一遍：数据位（位 3、5、6、7）依次为 1、0、1、1。$p_1$ 覆盖位 1、3、5、7，其中数据 1、0、1 共两个 1，$p_1=0$；$p_2$ 覆盖位 2、3、6、7，数据 1、1、1 共三个 1，$p_2=1$；$p_4$ 覆盖位 4、5、6、7，数据 0、1、1 共两个 1，$p_4=0$。发出码字（位 1 至 7）为 0 1 1 0 0 1 1。若位 6 在途中翻成 0：$p_1$ 组（0,1,0,1）共两个 1，通过；$p_2$ 组（1,1,0,1）共三个 1，失败；$p_4$ 组（0,0,0,1）共一个 1，失败。校正子按 $p_4p_2p_1$ 读出 $(1,1,0)=6$，直接指向位 6，翻转回来即完成纠正。

实用的 ECC 内存把同一思想做到 64 位数据加 8 位校验、72 位宽，实现单错纠正、双错检测（SEC-DED）。贵在三处：阵列宽了 1/8，每条内存要多挂芯片；控制器对每次读做译码纠正，对部分写先读-改-写再重算校验，多一层逻辑和时延；平台还要为纠错路径做完整验证。对长时间运行、数据不能错的机器这是便宜的保险，桌面机则普遍省掉——这也是普通内存条与 ECC 内存条不能随意混插的原因。

\section{二、地址译码定义整机地图}

CPU 输出的是地址数值，整机必须规定这个数值由哪个目标响应。这个规定叫地址映射。一个简单教学板可以把低地址映射到 ROM，中间映射到 RAM，高地址小窗口映射到 GPIO、定时器和串口等设备寄存器。地址译码器把高位地址和访问类型变成片选信号。

\begin{lstlisting}
example 16-bit address map:
  0x0000-0x7FFF  boot ROM
  0x8000-0xEFFF  SRAM
  0xF000         GPIO_OUT
  0xF004         GPIO_IN
  0xF010         TIMER_COUNT
  0xF014         TIMER_CTRL
\end{lstlisting}

\begin{pdfdiagram}
  \node[hw state, minimum width=2.0cm] (addr) at (0,0) {地址 A[15:0]};
  \node[hw compute, minimum width=2.15cm] (dec) at (3.0,0) {高位译码};
  \node[hw state, minimum width=1.6cm] (rom) at (6.0,1.05) {ROM CS};
  \node[hw state, minimum width=1.6cm] (ram) at (6.0,0) {RAM CS};
  \node[hw control, minimum width=1.6cm] (io) at (6.0,-1.05) {IO CS};
  \draw[hw bus] (addr) -- (dec);
  \draw[hw bus] (dec.east) |- (rom.west) node[pos=0.3,above,hw label]{A15=0};
  \draw[hw bus] (dec) -- (ram) node[pos=0.55,above,hw label]{1xxx（非 1111）};
  \draw[hw bus] (dec.east) |- (io.west) node[pos=0.3,below,hw label]{1111};
  \node[hw label, text width=7.8cm] at (3.4,-1.95) {设计目标：任意一个地址最多选中一个目标；未映射区域也要有错误、固定值或超时规则。};
\end{pdfdiagram}
\diagramnote{地址译码不是软件表格，而是片选电路。片选重叠会造成总线争用或误写设备。}

\begin{longtable}{L{0.18\linewidth}L{0.38\linewidth}L{0.34\linewidth}}
\toprule
映射对象 & 响应规则 & 典型错误 \\
\midrule
ROM & 复位入口和只读代码地址命中 & 复位后 CPU 取不到第一条指令 \\
RAM & 普通数据读写地址命中 & 写使能误打到 ROM 或设备窗口 \\
GPIO & 设备寄存器地址命中并产生副作用 & 地址译码过宽，普通内存写误触发外设 \\
未映射区 & 返回错误、超时或固定空值 & 访问悬空，CPU 永久等待或读到随机总线值 \\
\bottomrule
\end{longtable}

\topic{片选表达式就是硬件约定}
若地址空间为 64 KiB，可以用高位地址产生粗粒度片选。例如 \code{A15=0} 选 ROM，\code{A15=1} 且 \code{A14,A13,A12} 不全为 1 选 RAM（\code{0x8000-0xEFFF}，28 KiB），\code{A15=1,A14=1,A13=1,A12=1} 选 I/O 窗口。无论用门电路、译码器、CPLD 还是 FPGA 实现，都必须证明任意合法周期最多只有一个目标片选有效。

\begin{lstlisting}
rom_cs = !A15
ram_cs = A15 && !(A14 && A13 && A12)
io_cs  = A15 && A14 && A13 && A12
\end{lstlisting}

这三个表达式看起来像逻辑练习，实际上就是整机安全边界。若 ROM 和 RAM 可同时片选，读周期可能出现两个芯片同时驱动数据线；若 RAM 和 GPIO 片选重叠，普通 STORE 可能既改内存又改外设输出；若未映射地址没有超时或错误返回，CPU 可能等待一个永远不会响应的目标。

\topic{时序预算决定是否需要等待状态}
地址译码需要时间，存储器访问需要时间，数据到达 CPU 输入端后还要满足建立时间。若这些时间总和超过时钟周期，就必须降低频率、换更快器件，或插入等待状态。等待状态不是软件延时，而是硬件事务延长：地址、片选、读写方向和数据驱动关系必须在等待期间保持可解释。

\begin{lstlisting}
read_budget =
  cpu_address_output_delay
  + address_decode_delay
  + memory_access_time
  + input_setup_time
  + board_margin
\end{lstlisting}

若 \code{read_budget} 大于一个周期，控制器就不能在下一拍采样数据。慢 ROM、慢外设和长板级走线都可能需要 READY/WAIT 机制。正式规格还应说明最大等待长度和超时行为，避免坏设备把整机永久挂住。

\begin{longtable}{L{0.2\linewidth}L{0.34\linewidth}L{0.36\linewidth}}
\toprule
处理方式 & 适合场景 & 代价 \\
\midrule
降低时钟 & 全系统都慢，教学板容易观察 & 所有操作都变慢 \\
更快器件 & SRAM 或译码器成为瓶颈 & 成本、功耗或电平兼容性变化 \\
插入等待 & 只有部分目标较慢 & 控制器必须保持事务状态 \\
超时返回错误 & 目标可能不存在或失效 & 需要错误状态和恢复入口 \\
\bottomrule
\end{longtable}

\topic{案例：用 74HC138 译出整机片选}
本节开头给出了 16-bit 地图：\code{0x0000-0x7FFF} 是 boot ROM（32 KiB，A15=0），\code{0x8000-0xEFFF} 是 SRAM（28 KiB），\code{0xF000} 起是 I/O 区（4 KiB）。现在给出这张地图的引脚级实现。最经典的器件是 3-8 译码器 74HC138：它有三个地址输入 C、B、A 和八个低有效输出 Y0--Y7，任一时刻只有一个输出为低。把 A[15:13] 依次接上 C、B、A，每个输出就恰好管辖一个 8 KiB 窗口，窗内偏移由 A[12:0] 决定；使能端 G1 接高电平、/G2A 与 /G2B 接低电平，译码器才工作——这两个使能端也常改接总线读写选通，让片选只在存储周期内出现。输出低有效带来一个方便的性质：把若干相邻窗口合并成一个片选只需一个与门，任何一个被选中的输出变低，与门输出即为低。于是 Y0--Y3 经四输入与门合成 ROM\_CS，正好覆盖 A15=0 的整个低半空间；Y4--Y6 经三输入与门合成 RAM\_CS；Y7 则送入二级译码，由 A12 把 \code{0xE000-0xFFFF} 再分成 SRAM 尾窗和 I/O 区，合起来正好复现上面这张地图。

\begin{pdfdiagram}
  % 地址源与 138 本体（引脚级）
  \node[hw state, minimum width=2.2cm] (addr) at (0.65,1.2) {地址 A[15:0]};
  \node[hw compute, minimum width=2.2cm, minimum height=4.9cm] (hc) at (3.7,-0.2) {74HC138};
  % A[15:13] 接 C/B/A，A12 走底部干线去二级译码
  \draw[hw bus] (addr.east) -- (2.6,1.2) node[pos=0.65,above=1pt,hw label,fill=white,inner sep=1pt]{A14};
  \draw[hw bus] (2.05,1.2) -- (2.05,1.7) -- (2.6,1.7) node[pos=0.6,above=1pt,hw label,fill=white,inner sep=1pt]{A15};
  \draw[hw bus] (2.05,1.2) -- (2.05,0.7) -- (2.6,0.7) node[pos=0.6,above=1pt,hw label,fill=white,inner sep=1pt]{A13};
  \draw[hw bus] (2.05,1.2) -- (2.05,-3.75) -- (6.65,-3.75) -- (6.65,-3.1) node[pos=0.8,above=1pt,hw label,fill=white,inner sep=1pt]{A12};
  \fill[BookMuted] (2.05,1.2) circle (0.05);
  \fill[BookMuted] (2.05,0.7) circle (0.05);
  \node[hw label, anchor=west] at (2.68,1.7) {C};
  \node[hw label, anchor=west] at (2.68,1.2) {B};
  \node[hw label, anchor=west] at (2.68,0.7) {A};
  % 使能引脚：G1=1，/G2A=/G2B=0
  \draw[hw return] (3.0,-3.05) -- (3.0,-2.65);
  \draw[hw return] (3.7,-3.05) -- (3.7,-2.65);
  \draw[hw return] (4.4,-3.05) -- (4.4,-2.65);
  \node[hw label, anchor=south] at (3.0,-2.62) {G1};
  \node[hw label, anchor=south] at (3.7,-2.62) {/G2A};
  \node[hw label, anchor=south] at (4.4,-2.62) {/G2B};
  \node[hw label, anchor=north] at (3.7,-3.12) {G1=1，/G2A=/G2B=0};
  % 输出引脚 Y0--Y7
  \node[hw label, anchor=east] at (4.72,1.925) {Y0};
  \node[hw label, anchor=east] at (4.72,1.375) {Y1};
  \node[hw label, anchor=east] at (4.72,0.825) {Y2};
  \node[hw label, anchor=east] at (4.72,0.275) {Y3};
  \node[hw label, anchor=east] at (4.72,-0.275) {Y4};
  \node[hw label, anchor=east] at (4.72,-0.825) {Y5};
  \node[hw label, anchor=east] at (4.72,-1.375) {Y6};
  \node[hw label, anchor=east] at (4.72,-1.925) {Y7};
  % 片选合并与二级译码
  \node[hw compute, minimum width=1.1cm, minimum height=1.9cm] (or0) at (6.65,1.1) {与门};
  \node[hw compute, minimum width=1.1cm, minimum height=1.65cm] (or1) at (6.65,-0.825) {与门};
  \node[hw state, minimum width=1.5cm] (romcs) at (8.85,1.1) {ROM\_CS};
  \node[hw state, minimum width=1.5cm] (ramcs) at (8.85,-0.825) {RAM\_CS};
  \node[hw compute, minimum width=1.9cm, minimum height=1.1cm] (iodec) at (6.65,-2.55) {二级译码};
  \node[hw state, minimum width=1.7cm] (tail) at (9.45,-2.3) {SRAM 尾窗};
  \node[hw state, minimum width=1.7cm] (iocs) at (9.45,-2.8) {I/O 区 CS};
  \draw[hw bus] (4.8,1.925) -- (6.1,1.925);
  \draw[hw bus] (4.8,1.375) -- (6.1,1.375);
  \draw[hw bus] (4.8,0.825) -- (6.1,0.825);
  \draw[hw bus] (4.8,0.275) -- (6.1,0.275);
  \draw[hw bus] (4.8,-0.275) -- (6.1,-0.275);
  \draw[hw bus] (4.8,-0.825) -- (6.1,-0.825);
  \draw[hw bus] (4.8,-1.375) -- (6.1,-1.375);
  \draw[hw bus] (4.8,-1.925) -- (5.3,-1.925) -- (5.3,-2.55) -- (5.7,-2.55);
  \draw[hw bus] (or0) -- (romcs);
  \draw[hw bus] (or1) -- (ramcs);
  \draw[hw bus] (7.6,-2.3) -- (8.6,-2.3) node[pos=0.5,above=1pt,hw label,fill=white,inner sep=1pt]{A12=0};
  \draw[hw bus] (7.6,-2.8) -- (8.6,-2.8) node[pos=0.5,above=1pt,hw label,fill=white,inner sep=1pt]{A12=1};
\end{pdfdiagram}
\diagramnote{引脚级地址译码。A[15:13] 经 74HC138 译成 8 个 8 KiB 窗口（输出低有效），与门把连续窗口合成 ROM\_CS 与 RAM\_CS，Y7 窗口由 A12 二级译码再分；使能端决定译码器何时工作。}

\begin{longtable}{L{0.09\linewidth}L{0.18\linewidth}L{0.19\linewidth}L{0.44\linewidth}}
\toprule
输出 & C B A（A[15:13]） & 地址区间 & 去向 \\
\midrule
Y0 & 000 & \code{0x0000-0x1FFF} & boot ROM 第 0 页 \\
Y1 & 001 & \code{0x2000-0x3FFF} & boot ROM 第 1 页 \\
Y2 & 010 & \code{0x4000-0x5FFF} & boot ROM 第 2 页 \\
Y3 & 011 & \code{0x6000-0x7FFF} & boot ROM 第 3 页 \\
Y4 & 100 & \code{0x8000-0x9FFF} & SRAM 第 0 页 \\
Y5 & 101 & \code{0xA000-0xBFFF} & SRAM 第 1 页 \\
Y6 & 110 & \code{0xC000-0xDFFF} & SRAM 第 2 页 \\
Y7 & 111 & \code{0xE000-0xFFFF} & 二级译码：A12=0 归 SRAM 尾窗（\code{0xE000-0xEFFF}），A12=1 为 I/O 区（\code{0xF000-0xFFFF}） \\
\bottomrule
\end{longtable}

这套译码仍然是部分译码：I/O 区里只有 A12 和少数低位参与选择，A[11:5] 等地址位没有进译码器。结果是同一个设备寄存器会在整个 4 KiB 窗口里反复出现——\code{0xF010} 的定时器寄存器，在 \code{0xF110}、\code{0xF210} 等只改变未译码位的地址上同样可见，这些重复地址称为别名。SRAM 一侧同理：若实际芯片容量小于所属窗口，未译码的高位也会让同一个存储单元出现在多个地址上。部分译码用更少的门和更短的延迟换来足够大的窗口，在小系统里是合理的工程选择；但它把一道边界推给了规格：别名是否允许必须写明。允许别名，软件就不得用未译码位区分设备，调试时也不能从“访问了哪个地址”反推“选中了哪个设备”；不允许别名，就必须补足低位译码，让未映射组合返回错误或固定值，而不是把一次写错的访问悄悄变成对某个真实设备的读写。

\section{三、MMIO 把设备寄存器放进地址空间}

设备寄存器可以像内存地址一样被 CPU 读写，但它们不是普通 RAM。读状态寄存器可能清除某些位，写命令寄存器可能启动外设动作，写数据寄存器可能进入 FIFO，读数据寄存器可能弹出一个元素。这类地址称为内存映射 I/O，简称 MMIO。MMIO 的关键不是地址长得像内存，而是每个寄存器都有副作用和顺序规则。

\begin{longtable}{L{0.18\linewidth}L{0.36\linewidth}L{0.36\linewidth}}
\toprule
设备寄存器 & 读写语义 & 关键问题 \\
\midrule
\code{DIR} & 设置 GPIO 引脚方向 & 输出引脚和输入引脚不会混用 \\
\code{OUT} & 写输出锁存器 & 写入只改变目标输出位 \\
\code{IN} & 读取同步后的外部输入 & 按键抖动和亚稳态不能直接进入 CPU \\
\code{STAT} & 返回 ready、busy、error、irq 等状态 & 清除规则不会丢掉同时到来的新事件 \\
\code{DATA} & 数据收发寄存器或 FIFO 端口 & 空读、满写和溢出有明确定义 \\
\bottomrule
\end{longtable}

\topic{MMIO 写响应不等于设备已经完成}
CPU 写一个命令寄存器，通常只表示设备已经接收命令。设备真正完成可能要等若干周期、若干微秒，甚至更久。完成可以通过状态位、轮询、DMA completion 或中断通知。若把 MMIO 写返回当作“设备动作完成”，软件会在设备尚未写完缓冲区、尚未刷新状态或尚未清除错误时继续执行，形成难复现的硬件错误。

\begin{lstlisting}
device command sequence:
  CPU writes COMMAND register
  bus write response returns
  device starts internal work
  device sets DONE or ERROR
  interrupt or polling tells CPU to inspect result
\end{lstlisting}

\begin{pdfdiagram}
  \node[hw state, minimum width=1.55cm] (cpu) at (0,0) {CPU};
  \node[hw compute, minimum width=1.75cm] (bus) at (2.3,0) {总线响应};
  \node[hw control, minimum width=1.95cm] (cmd) at (4.85,0) {命令寄存器};
  \node[hw control, minimum width=1.95cm] (dev) at (7.45,0) {设备状态机};
  \node[hw state, minimum width=1.5cm] (done) at (9.75,0) {DONE};
  \draw[hw bus] (cpu) -- node[above=2pt,hw label,fill=white,inner sep=1pt]{写命令} (bus);
  \draw[hw bus] (bus) -- node[above=2pt,hw label,fill=white,inner sep=1pt]{写已接受} (cmd);
  \draw[hw bus] (cmd) -- node[above=2pt,hw label,fill=white,inner sep=1pt]{开始工作} (dev);
  \draw[hw return] (dev) -- node[above=2pt,hw label,fill=white,inner sep=1pt]{完成/错误} (done);
  \node[hw label, text width=9.0cm] at (4.7,-1.0) {总线写响应只证明命令到达设备接口；真正完成必须看状态位、中断或 DMA completion。};
\end{pdfdiagram}
\diagramnote{MMIO 写命令的两层完成边界。总线完成不等于设备动作完成。}

\topic{端口 I/O 与 MMIO 是两种地址约定}
8086 家族保留了独立 I/O 地址空间，\code{IN} 和 \code{OUT} 指令访问端口；许多 RISC 和单片机系统则把设备寄存器映射进普通地址空间，用 LOAD/STORE 访问。两者都能控制设备，但硬件约定不同。端口 I/O 需要 I/O 周期和端口地址译码；MMIO 需要内存地址译码和内存属性规则，例如不可缓存、不可合并或必须保持顺序。

\topic{外部输入必须先调理和同步}
按键、传感器、串口线和外部中断引脚都不按 CPU 时钟整齐变化。按键还会抖动，异步输入还可能让触发器进入亚稳态。教学板上的 GPIO 输入至少应经过电平保护、上拉/下拉、去抖策略和同步触发器，再进入设备寄存器。否则 CPU 读到的不是稳定硬件事实，而是未定义的边界行为。

\begin{warningbox}
不要把设备寄存器当作普通变量。普通 RAM 的读写目标是保存数据；MMIO 的读写目标是驱动硬件状态机。编译器、cache、写缓冲和乱序执行都可能改变普通内存访问的时机，但设备访问通常需要不可缓存属性、顺序约束和明确的完成检查。
\end{warningbox}

\topic{从 GPIO、定时器和 UART 看设备寄存器}
设备寄存器最好用具体布局理解。一个 GPIO 控制器可以有输出寄存器、输入寄存器、方向寄存器和中断状态寄存器；定时器可以有计数值、比较值、控制位和中断清除位；UART 可以有发送数据、接收数据、状态和波特率配置寄存器。CPU 访问这些寄存器时走的是地址事务，但设备内部看到的是控制状态机输入。

\begin{longtable}{L{0.18\linewidth}L{0.36\linewidth}L{0.36\linewidth}}
\toprule
设备 & 寄存器边界 & 常见问题 \\
\midrule
GPIO & \code{DIR/OUT/IN/IRQ_STAT} & 输入是否同步，输出是否只改目标位 \\
定时器 & \code{COUNT/CMP/CTRL/IRQ_CLEAR} & 到期事件是否锁存，清除是否丢新事件 \\
UART & \code{TX/RX/STAT/BAUD} & 空读、满写、错误位和中断顺序是否清楚 \\
\bottomrule
\end{longtable}

\topic{UART 帧格式与波特率约定}
UART 没有时钟线，收发双方只靠“每一位持续同样长的时间”这条约定保持对齐。线路空闲时停在高电平；发送方先用一个位时间的低电平打出起始位，再把 8 个数据位按 LSB 在前依次驱动上线，最后用至少一个位时间的高电平作为停止位。接收端用 16 倍于波特率的本地时钟对线路过采样：检测到下降沿后数 8 个采样周期，在起始位中点确认它仍是低电平，随后每隔 16 个采样周期在各位中点读取一次。中点采样让判决位置离前后边沿各半位，这正是异步串行能够容忍时钟误差的来源。

\begin{waveform}[x=0.85cm,y=0.95cm]
  % 一帧：空闲高 -> 起始位 0 -> D0..D7（0x55，LSB 在前）-> 停止位 1
  \draw[BookAccent, line width=0.9pt] (0,2.0) -- (1,2.0) -- (1,1.2) -- (2,1.2) -- (2,2.0) -- (3,2.0) -- (3,1.2) -- (4,1.2) -- (4,2.0) -- (5,2.0) -- (5,1.2) -- (6,1.2) -- (6,2.0) -- (7,2.0) -- (7,1.2) -- (8,1.2) -- (8,2.0) -- (9,2.0) -- (9,1.2) -- (10,1.2) -- (10,2.0) -- (11.8,2.0);
  \node[hw label, anchor=east] at (-0.2,1.6) {RXD};
  \node[hw label] at (0.5,2.25) {空闲};
  \node[hw label] at (1.5,2.25) {起始};
  \node[hw label] at (2.5,2.25) {D0};
  \node[hw label] at (3.5,2.25) {D1};
  \node[hw label] at (4.5,2.25) {D2};
  \node[hw label] at (5.5,2.25) {D3};
  \node[hw label] at (6.5,2.25) {D4};
  \node[hw label] at (7.5,2.25) {D5};
  \node[hw label] at (8.5,2.25) {D6};
  \node[hw label] at (9.5,2.25) {D7};
  \node[hw label] at (10.5,2.25) {停止};
  \draw[hw return] (1.5,0.75) -- (10.5,0.75);
  \foreach \x in {1.5,2.5,3.5,4.5,5.5,6.5,7.5,8.5,9.5,10.5} \draw[hw return] (\x,0.75) -- (\x,1.12);
  \node[hw label] at (6.0,0.28) {采样点（$16\times$ 过采样，位中点）};
\end{waveform}
\diagramnote{UART 一帧（数据 \code{0x55}，LSB 在前）。空闲为高电平；起始位的下降沿给出帧起点，数据位在中点被逐位采样，停止位把线路带回高电平。}

以 115200 baud 为例，每位宽度 $1/115200\approx 8.68\,\mu\text{s}$。接收端在起始位边沿对齐之后，最后一位（停止位）的采样点距该边沿 9.5 个位时间；若收发时钟的相对误差为 $\varepsilon$，这个采样点会漂移 $9.5\varepsilon$ 位，漂移超过半位即采到相邻位，所以理论上限约为 $0.5/9.5\approx 5\%$。工程上还要从中扣除边沿检测的量化误差（16 倍过采样下最多 $1/16$ 位）与电平翻转时间，留下的常用要求是每一侧波特率误差不超过约 2\%，或者收发双方干脆从同一时钟源分频。误差超限时，出错的往往不是帧首而是帧尾——越早的位漂移越小，最后的停止位最先越界，表现为间歇性的帧错误或数据错位。

\topic{定时器的预分频、计数与比较匹配}
定时器做的事只有一件：把固定频率的时钟变成软件可编程的周期间隔。它的数据路径很简单——预分频器先把输入时钟除以 $N$，递增计数器 CNT 在每个分频后时钟上加一，比较器持续比较 CNT 与比较寄存器 CMP；当 CNT 数到 CMP，比较器输出 match。match 同时驱动两件事：置起事件状态位（若控制寄存器 CR 中的中断使能有效，同时拉起 IRQ），以及在 auto-reload 模式下把 CNT 清零，让计数立刻开始下一轮。于是中断周期严格等于（预分频 $\times$ (CMP+1)）个输入时钟——CMP 要加一，因为计数从 0 开始，数到 CMP 共经过 CMP+1 拍。以 1 MHz 输入、预分频 8 为例：分频后的计数时钟是 125 kHz，若 CMP=999，周期为 $8\times(999+1)=8000$ 个输入时钟，即 8 ms（125 Hz）；要得到 1 kHz 中断，应令 CMP=124，此时 $8\times(124+1)=1000$ 个时钟，正好 1 ms。事件状态位由软件按设备规则清除，清除顺序属于§四要讨论的中断约定。

\begin{pdfdiagram}
  \node[hw control, minimum width=1.6cm] (clk) at (0.3,0) {时钟 CLK};
  \node[hw compute, minimum width=1.9cm] (pre) at (2.5,0) {预分频器 $\div N$};
  \node[hw state, minimum width=2.2cm] (cnt) at (5.15,0) {递增计数器 CNT};
  \node[hw compute, minimum width=1.5cm] (cmp) at (7.5,0) {比较器};
  \node[hw state, minimum width=1.9cm] (cmpreg) at (7.5,-1.05) {CMP 寄存器};
  \node[hw state, minimum width=1.7cm] (evt) at (9.9,0) {事件状态位};
  \node[hw control, minimum width=1.9cm] (irq) at (9.9,-1.05) {中断请求 IRQ};
  \node[hw control, minimum width=1.9cm] (cr) at (3.4,-1.4) {控制寄存器 CR};
  \draw[hw bus] (clk) -- (pre);
  \draw[hw bus] (pre) -- (cnt);
  \draw[hw bus] (cnt) -- (cmp);
  \draw[hw bus] (cmpreg.north) -- (cmp.south);
  \draw[hw bus] (cmp) -- node[above=2pt,hw label,fill=white,inner sep=1pt]{match} (evt);
  \draw[hw return] (evt.south) -- (irq.north);
  \draw[hw return] (cmp.north) -- (7.5,0.95) -- (5.15,0.95) node[pos=0.5,above=1pt,hw label,fill=white,inner sep=1pt]{归零重载（auto-reload）} -- (cnt.north);
  \draw[hw return] (cr.north) -- (3.4,-0.65) -- (2.5,-0.65) -- (pre.south);
  \draw[hw return] (3.4,-0.65) -- (5.15,-0.65) node[pos=0.5,below=2pt,hw label,fill=white,inner sep=1pt]{使能/重装} -- (5.15,-0.4);
  \fill[BookTeal] (3.4,-0.65) circle (0.05);
\end{pdfdiagram}
\diagramnote{定时器数据路径：时钟经预分频驱动递增计数器，比较器在 CNT 等于 CMP 时产生 match，一路置事件状态位并按 CR 的中断使能拉起 IRQ，一路按 auto-reload 把 CNT 归零，开始下一周期。中断周期 = 预分频 $\times$ (CMP+1) 个输入时钟。}

\topic{I2C 用两根开漏线支撑一条多设备总线}
I2C 只有两根信号线：串行数据线 SDA 和串行时钟线 SCL。两根线都是开漏输出加上拉电阻——任何设备都只能把线拉低，不能主动驱动为高，线电平相当于所有设备输出的“线与”。空闲时两线靠上拉电阻停在高电平；主设备产生全部时钟，数据位只允许在 SCL 为低时改变，在 SCL 为高时被采样。事务以起始条件开头——SCL 保持高电平时 SDA 从高跳低；以停止条件结尾——SCL 保持高电平时 SDA 从低跳高。每个字节由 8 个数据位加 1 个应答位组成：第 9 个时钟里发送方释放 SDA，接收方把 SDA 拉低表示 ACK，保持高表示 NACK。寻址用 7 位设备地址，与 1 位读写方向合成地址阶段的一个字节。开漏结构正是这条总线能够共享的原因：多主同时发送时，发 0 的一方自然赢过发 1 的一方，竞争不需要额外仲裁线；从设备还可以在字节之间把 SCL 按住不放，强制主设备等待（clock stretching），慢设备因此不会丢数据。

UART 是点对点链路，没有地址概念；I2C 把多个设备挂在同一对线上，用起始条件和地址字节决定本次事务跟谁说话。下表对照二者的边界。

\begin{longtable}{L{0.16\linewidth}L{0.37\linewidth}L{0.37\linewidth}}
\toprule
方面 & I2C & UART \\
\midrule
拓扑 & 一条共享总线挂多个主从设备 & 点对点，一发一收 \\
时钟 & 主设备产生 SCL，收发同步于时钟沿 & 无时钟线，靠双方约定的波特率 \\
寻址 & 起始条件后发送 7 位地址加读写位 & 无地址，线路上只有一个对端 \\
确认 & 每字节后跟 ACK/NACK，主设备能发现无应答 & 无应答机制，只能靠上层协议 \\
线数与速度 & 2 线；标准 100 kHz、快速 400 kHz & 2 线收发交叉；常见 115200 baud \\
\bottomrule
\end{longtable}

一次典型的“主发地址、寄存器号、再读回数据”事务如下，传感器和 EEPROM 的随机读几乎都是这个形状：

\begin{lstlisting}
I2C register read transaction:
  START
  master sends [7-bit address + W]  -> slave ACK
  master sends [register number]    -> slave ACK
  repeated START
  master sends [7-bit address + R]  -> slave ACK
  slave sends  [data byte]          -> master NACK
  STOP
\end{lstlisting}

中间的重复起始（repeated START）把“先告诉对方读哪个寄存器”和“再开始读”合并成一次不间断事务，避免其他主设备在两段之间插入；读方向下最后一个字节由主设备回 NACK 再发停止条件，通知从设备释放 SDA、结束发送。这次事务的开销可以直接算：一帧 9 位，共 4 帧 36 位，未计起始和停止条件的少量时间；标准模式 100 kHz 下每位 $10\,\mu\text{s}$，合计约 $36\times10=360\,\mu\text{s}$，快速模式 400 kHz 下约 $90\,\mu\text{s}$。也就是说，主设备每 0.36 ms 才能完成一次这样的寄存器读取，这正是 I2C 传感器数据更新率的量级来源；I2C 的意义不在于快，而在于只用两根线就把一整板低速设备纳入同一个地址空间。

\topic{SPI 用片选换速度，用移位环换全双工}
SPI 是四线同步串行接口：SCLK 时钟、MOSI 主出从入、MISO 主入从出、\code{CS} 片选（低有效）。总线上没有地址阶段——主设备想跟谁通信就把谁的片选拉低，同一时刻只允许被选中的从设备驱动 MISO。数据通路是一对首尾相接的移位寄存器：每来一个时钟沿，主设备从 MOSI 移出一位，同时从 MISO 移入一位，收发在物理上永远同时进行，全双工是这个结构的自然结果而不是协议选项。时钟行为由两个参数规定：CPOL 决定空闲时 SCLK 停在高还是低，CPHA 决定在第几个时钟沿采样数据，组合出四种模式。通信双方必须事先约定同一种模式——这条协商写在芯片手册里，不在总线上进行。

\begin{longtable}{L{0.12\linewidth}L{0.12\linewidth}L{0.12\linewidth}L{0.24\linewidth}L{0.28\linewidth}}
\toprule
模式 & CPOL & CPHA & 空闲电平 & 数据采样沿 \\
\midrule
0 & 0 & 0 & 低 & 第一个沿（上升沿） \\
1 & 0 & 1 & 低 & 第二个沿（下降沿） \\
2 & 1 & 0 & 高 & 第一个沿（下降沿） \\
3 & 1 & 1 & 高 & 第二个沿（上升沿） \\
\bottomrule
\end{longtable}

以最常用的模式 0 读一个字节为例：主设备先拉低 \code{CS}；从设备把数据的最高位驱动到 MISO（发送方在时钟下降沿换位）；主设备发出 8 个时钟，在每个上升沿采样 MISO，依次得到 D7 到 D0；与此同时主设备把 MOSI 上的 8 位移给从设备，纯读时通常发全 0 占位字节；8 拍结束后主设备拉高 \code{CS}，事务结束。

\begin{lstlisting}
SPI mode-0 byte read (MSB first):
  CS low
  for bit = D7 downto D0:
      slave changes MISO on falling edge of SCLK
      master samples MISO on rising edge of SCLK
      master shifts one dummy bit out on MOSI
  CS high
\end{lstlisting}

SPI 没有 ACK，也没有起始停止开销，8 个时钟就是 8 位；代价写在引脚数量上，每多一个从设备多一根片选。与 I2C 对照，两者在速度、引脚和协议复杂度上各站一端：

\begin{longtable}{L{0.20\linewidth}L{0.35\linewidth}L{0.35\linewidth}}
\toprule
方面 & SPI & I2C \\
\midrule
速度 & 常见数 MHz 到数十 MHz & 标准 100 kHz，快速 400 kHz \\
引脚 & 4 线，且每个从设备一根片选 & 2 线，设备数量不占引脚 \\
寻址 & 无地址，靠片选 & 7 位地址加读写位 \\
确认与等待 & 无 ACK，无从设备等待机制 & 每字节 ACK/NACK，从设备可拉住 SCL 等待 \\
协议复杂度 & 几乎只有移位，模式须事先约定 & 起始、停止、重复起始、应答状态机 \\
\bottomrule
\end{longtable}

速度差距可以直接算：SCLK 取 10 MHz 时位时间 $100\,\text{ns}$，读一个字节恰需 8 拍 $=0.8\,\mu\text{s}$；I2C 快速模式 400 kHz 下位时间 $2.5\,\mu\text{s}$，一个字节加应答共 9 位 $=22.5\,\mu\text{s}$，相差约 28 倍。Flash、显示屏、ADC 这类数据量较大的从设备多用 SPI，原因就在这里；而传感器这类低速、多节点、布线路径长的场合，I2C 的两线共享更省板级资源。可见两条总线没有优劣之分，只有“用引脚和约定换什么”的取舍不同。

\topic{USB 用枚举和描述符换来“通用”二字}
UART、SPI、I2C 的共同前提是主设备事先知道对端是什么；USB 把这个前提取消了。总线严格以主机为中心，设备不能主动发起任何事务。插入检测靠电气约定：设备端在 D+ 或 D- 上接上拉电阻，全速设备拉 D+，低速设备拉 D-，集线器看到哪条线被拉高就知道来了什么速度的设备（高速设备先以全速现身，在复位期间再协商提速），并把端口状态变化上报主机。随后主机执行枚举：先复位该端口，设备以默认地址 0 应答；主机通过默认控制端点 0 读取设备描述符，再发 Set Address 给设备分配 1 到 127 之间的新地址；之后按新地址继续读取配置描述符、接口描述符、端点描述符和字符串描述符。描述符是一套标准格式的自描述数据：厂商 ID、产品 ID、设备类、需要多少端点、每个端点的类型和最大包长，操作系统据此匹配驱动。插拔能自动识别，不是因为总线聪明，而是因为每台设备都能用同一种格式回答“我是谁、我要什么”；“通用”二字，正来自枚举和描述符这套约定。

枚举完成后，数据按端点组织。端点是设备内部的单向通道，端点 0 固定留给控制传输，其余端点在描述符中声明自己的传输类型：

\begin{longtable}{L{0.16\linewidth}L{0.40\linewidth}L{0.34\linewidth}}
\toprule
传输类型 & 语义 & 典型设备 \\
\midrule
控制 & 有应答、用于配置和命令，端点 0 专用 & 所有设备的枚举阶段 \\
中断 & 主机按声明周期轮询，小数据量，出错重试 & 键盘、鼠标 \\
批量 & 利用剩余带宽，无周期保证，出错重传 & U 盘、打印机 \\
等时 & 每帧预留带宽，不重传，容忍少量错误 & 音频设备、摄像头 \\
\bottomrule
\end{longtable}

把四种总线放在一起，复杂度阶梯很清楚：

\begin{longtable}{L{0.12\linewidth}L{0.18\linewidth}L{0.20\linewidth}L{0.16\linewidth}L{0.24\linewidth}}
\toprule
总线 & 线数 & 寻址方式 & 设备自描述 & 典型定位 \\
\midrule
UART & 2 & 无，点对点 & 无 & 调试口、模块间低速链路 \\
SPI & 4，每从设备加 1 片选 & 片选 & 无 & 板内高速从设备 \\
I2C & 2 & 7 位地址 & 无 & 板内多节点低速设备 \\
USB & 4（含电源和地） & 枚举分配 7 位地址 & 标准格式描述符 & 机箱外可热插拔外设 \\
\bottomrule
\end{longtable}

带宽数字能说明“通用”的另一面：高速模式 480 Mb/s 折合每微帧 $125\,\mu\text{s}$ 约 7500 B，批量端点最大包 512 B，一个微帧最多容纳 13 个整包，理论上限约 $13\times512\,\text{B}\div125\,\mu\text{s}\approx53\,\text{MB/s}$。同一条总线既要服务每 $10\,\text{ms}$ 被轮询一次、一次只发几个字节的键盘，又要服务成批搬数据的 U 盘，靠的正是把调度规则写进四种传输类型，而不是为每类设备各做一种接口。

\section{四、中断和 DMA 把外部事件接回 CPU}

轮询可以让 CPU 一直读状态寄存器，直到设备 ready。但若设备很慢，CPU 会浪费大量周期。中断提供另一种路径：设备在状态变化后请求中断，中断控制器汇总、屏蔽和排序请求，CPU 在可接受边界保存当前执行上下文，跳到中断入口，服务程序读取设备状态并清除事件。

\begin{longtable}{L{0.18\linewidth}L{0.38\linewidth}L{0.34\linewidth}}
\toprule
对象 & 硬件职责 & 关键问题 \\
\midrule
设备 & 产生完成、错误或数据到达事件 & 状态位和中断请求是否同步一致 \\
中断控制器 & 汇总、屏蔽、优先级和投递 & 屏蔽位、优先级、EOI 顺序是否正确 \\
CPU 入口 & 保存边界状态并跳到入口 & PC、标志和特权边界是否可恢复 \\
服务程序 & 读状态、搬数据、清除事件 & 不丢新事件，不重复进入同一旧事件 \\
\bottomrule
\end{longtable}

\topic{中断链路从设备事件一直接到服务程序返回}中断把“CPU 主动轮询”反转为“设备主动请求”。完整链路是：事件先在设备内锁存成状态位并拉起 IRQ；中断控制器按屏蔽位和优先级决定是否投递；CPU 只在指令边界采样 INT，接收后保存 PC/Flags，查向量表跳到服务程序；服务程序读状态、处理数据、按设备规则清事件，最后发 EOI 并恢复现场返回。链路上任何一环顺序错误，都会表现为丢中断或同一事件重复进入。

\begin{pdfdiagram}
  \node[hw memory, minimum width=1.7cm] (dev) at (0,1.7) {设备};
  \node[hw state, minimum width=2.0cm] (latch) at (2.7,1.7) {事件锁存\\状态位};
  \node[hw control, minimum width=2.2cm] (pic) at (5.6,1.7) {中断控制器\\屏蔽/优先级};
  \node[hw state, minimum width=1.6cm] (cpu) at (8.7,1.7) {CPU};
  \draw[hw bus] (dev) -- node[above,hw label]{事件} (latch);
  \draw[hw bus] (latch) -- node[above,hw label]{IRQ} (pic);
  \draw[hw bus] (pic) -- node[above,hw label]{INT} (cpu);
  \node[hw state, minimum width=2.4cm] (save) at (8.7,0.0) {保存 PC/Flags};
  \node[hw compute, minimum width=2.2cm] (vec) at (5.6,0.0) {查向量表跳转};
  \node[hw compute, minimum width=3.2cm] (isr) at (2.4,0.0) {服务程序：读状态、\\处理、清事件、发 EOI};
  \draw[hw bus] (cpu.south) -- node[right,hw label]{指令边界接收} (save.north);
  \draw[hw bus] (save) -- (vec);
  \draw[hw bus] (vec) -- (isr);
  \draw[hw return] (isr.north) -- (2.4,0.85) -- (5.6,0.85) node[above, hw label, pos=0.5]{EOI} -- (pic.south);
  \draw[hw return] (isr.south) -- (2.4,-1.05) -- (10.1,-1.05) node[above, hw label, pos=0.5]{恢复现场，返回断点} -- (10.1,1.7) -- (cpu.east);
\end{pdfdiagram}
\diagramnote{中断的完整链路。实线是请求与服务路径，虚线是返回路径：EOI 让中断控制器恢复投递，恢复现场让 CPU 回到被打断的位置。服务程序必须按设备规定的顺序清事件——先清设备还是先 EOI，是每种芯片数据手册都会写明的细节。}

\topic{中断清除顺序是硬件约定}
若服务程序先清设备状态，再清中断控制器，或先通知中断控制器再清设备，效果可能不同。某些设备要求先读状态后读数据，某些要求写 1 清位，某些要求读数据寄存器自动清除 ready。教材级规格必须写清楚事件锁存、状态读取、清除、EOI 和重新开放中断的顺序。否则系统在低速测试中正常，在高频事件下丢中断或重复中断。

\topic{中断控制器用三张寄存器决定谁先被服务}
设备多起来之后，“谁先被服务”不能靠先到先得。以 8259 可编程中断控制器为例，它用三张 8 位寄存器管理 8 条中断线：IRR（中断请求寄存器）锁存已到但尚未投递的请求；IMR（中断屏蔽寄存器）按位禁止对应线路参与竞争；ISR（在服务寄存器）记录当前正在被服务的请求。优先级裁决器在未屏蔽的 IRR 位中选出最高优先级的一位，与 ISR 中已在服务的最高位比较，只有新请求优先级更高时才向 CPU 拉起 INT。默认是固定优先级——IR0 最高、IR7 最低；控制器也可配置成轮转优先级，刚被服务的线路自动排到队尾，避免某条热线路长期独占。

\begin{longtable}{L{0.14\linewidth}L{0.40\linewidth}L{0.36\linewidth}}
\toprule
寄存器 & 记录内容 & 关键问题 \\
\midrule
IRR & 已锁存、待投递的请求位 & 请求是否真正被锁存，边沿触发与电平触发是否分清 \\
IMR & 每条线路的屏蔽位 & 屏蔽期间到来的请求是挂起还是丢失 \\
ISR & 正在服务的请求位 & EOI 之前同级和低级请求一律等待 \\
\bottomrule
\end{longtable}

EOI 是这套机制里最要紧的动作。服务程序发出 EOI，控制器才清除对应的 ISR 位；ISR 位不清，同级和更低优先级的请求即使进了 IRR 也只能等待——优先级秩序正是靠 ISR 位维持的。若希望高级中断能打断正在运行的低级服务程序，即允许嵌套，低级服务程序必须在保存好现场之后重新打开 CPU 的中断允许位。嵌套时保存的内容必须完整，因为返回时要逐层恢复：CPU 接收中断时已把返回地址和标志寄存器压栈，服务程序还要保存自己用到的全部通用寄存器。每一层嵌套占用一帧栈，层数越深栈消耗越大，这是嵌套中断容易忽略的隐性成本。

\begin{lstlisting}
isr_low:
    push rax            ; save registers this handler uses
    push rbx
    sti                 ; allow higher-priority nesting only now
    ; ... read device, handle data, clear device event ...
    cli                 ; block nesting on the return path
    mov al, 0x20        ; non-specific EOI command
    out 0x20, al        ; write 8259 command port, clears top ISR bit
    pop rbx
    pop rax
    iretq               ; restore flags and return address
\end{lstlisting}

优先级机制带来的风险是反转与饿死。反转的典型情形：低级服务程序长时间关中断执行，高级请求虽已进 IRR 却无法投递，实际等待时间反而更长。可以量化：115200 baud 下串口一个字节连起始、停止位共 10 位，到达间隔约 $86.8\,\mu\text{s}$；接收缓冲只有一字节深时，若某个低级服务程序关中断超过这个时间，串口数据就会被下一字节覆盖，表现为溢出错误——优先级更高的串口实际上输给了优先级更低的服务程序。饿死的典型情形则相反：固定优先级下高级中断密集到来时，低级请求在 IRR 里长期得不到服务，轮转优先级正是为这种情况准备的。工程上的通用对策有两条：服务程序尽量短，把重活推迟到中断返回之后；必须关中断的窗口要给出时间上界，而不是随手一关了事。

DMA 解决的是另一类浪费。若磁盘、网卡、显示控制器或音频设备需要搬运大量数据，CPU 不应逐字从设备读出再逐字写入内存。DMA 让设备或 DMA 控制器成为总线发起者，按描述符或寄存器配置直接读写主存。CPU 负责准备缓冲区和描述符，设备负责搬运，完成后用状态位或中断通知 CPU。

\begin{lstlisting}
DMA receive sketch:
  CPU allocates buffer
  CPU writes descriptor address and length
  CPU ensures descriptor is visible to device
  device reads descriptor and writes buffer
  device posts completion
  CPU handles interrupt and inspects buffer
\end{lstlisting}

\topic{DMA 控制器按通道组织，每个通道就是地址加计数}
8237 是理解 DMA 控制器的经典样本。它提供 4 个独立通道，每个通道对应一条设备请求线 DREQ；通道的核心状态只有两类 16 位寄存器——当前地址寄存器指出下一个要访问的内存地址，每传一字节自动加一，当前计数寄存器记录还剩多少字节，每传一字节自动减一。计数写入值比真实字节数少一：写 0 表示传 1 字节，写 \code{0xFFFF} 表示传 65536 字节，计数从 0 再减一回绕时产生终止计数 TC。模式寄存器决定传输方向（设备到内存或内存到设备）和传输方式。由于 16 位地址最多覆盖 64 KiB，PC 还给每个通道配了一个外部页寄存器，锁存物理地址的高 8 位，合成 24 位地址；由此带来一条经典限制：一次传输中低 16 位地址在页内推进，页寄存器保持不变，缓冲区不得跨越 64 KiB 页边界。

\begin{longtable}{L{0.20\linewidth}L{0.36\linewidth}L{0.34\linewidth}}
\toprule
寄存器组 & 记录内容 & 关键问题 \\
\midrule
地址寄存器 & 下一个内存单元的低 16 位地址 & 初值是否指向缓冲区起点，方向是否为递增 \\
计数寄存器 & 剩余字节数，写入值等于实际字节数减 1 & 忘记减 1 会多传一个字节，覆盖缓冲区之后的内容 \\
模式/命令 & 传输方向、单字/块/请求方式、通道号 & 方向写反会把设备数据当内存数据搬 \\
页寄存器 & 物理地址高 8 位（外部锁存） & 缓冲区是否跨 64 KiB 页边界 \\
屏蔽/请求 & 通道使能与软件请求 & 配置期间通道必须屏蔽，配完再开放 \\
\bottomrule
\end{longtable}

传输方式回答“拿到总线后干多久”。单字（single）模式每传一个字节就把总线还给 CPU，等设备下一次 DREQ 再申请；块（block）模式一旦拿到总线就连传到计数结束，其间 CPU 无法使用总线；请求（demand）模式随 DREQ 电平走走停停。慢设备配单字模式时，DMA 只是零星借用本属于 CPU 的总线周期，这称为周期窃取（cycle steal）：CPU 不被告知、不切换状态，只在恰好要用总线的那一拍多等一拍。块模式是另一个极端——它把 CPU 完全关在总线之外，换取最短的总搬运时间。

以软盘控制器常用的通道 2 为例，把 8192 字节从设备搬到物理地址 \code{0x23450} 的完整配置序列如下。页号 \code{0x02}、页内偏移 \code{0x3450}，结束地址 \code{0x2544F}，不跨页边界；计数写 $8192-1=8191=$ \code{0x1FFF}。拼接关系与这个缓冲区在页内的位置如图所示。

\begin{pdfdiagram}
  \node[hw state, minimum width=2.6cm] (pg) at (0,0.7) {页寄存器（外部锁存）\\A23--A16 = \code{0x02}};
  \node[hw state, minimum width=2.6cm] (ad) at (0,-0.7) {通道地址寄存器\\A15--A0 = \code{0x3450}};
  \node[hw compute, minimum width=1.2cm] (cat) at (3.6,0) {拼接};
  \node[hw state, minimum width=2.4cm] (pa) at (6.3,0) {物理地址\\\code{0x23450}};
  \draw[hw bus] (pg.east) -| node[pos=0.25,above=1pt,hw label]{高 8 位} (cat.north);
  \draw[hw bus] (ad.east) -| node[pos=0.25,below=1pt,hw label]{低 16 位} (cat.south);
  \draw[hw bus] (cat) -- (pa);
  % 页内位置：页 0x02 覆盖 0x20000--0x2FFFF
  \draw[BookRule, line width=0.5pt] (-1.5,-2.6) rectangle (6.8,-1.9);
  \draw[BookTeal, line width=0.7pt, fill=BookTeal!18] (0.22,-2.6) rectangle (1.23,-1.9);
  \node[hw label, anchor=west] at (1.45,-2.25) {缓冲区 \code{0x23450}--\code{0x2544F}（8192 B）};
  \node[hw label, anchor=south west] at (-1.5,-1.88) {页 \code{0x02}（64 KiB）};
  \node[hw label, anchor=south east] at (6.8,-1.88) {页边界禁跨};
  \node[hw label, anchor=north] at (-1.5,-2.68) {\code{0x20000}};
  \node[hw label, anchor=north] at (6.8,-2.68) {\code{0x2FFFF}};
  \draw[hw return] (pa.south) -- (6.3,-1.5) -- node[pos=0.5,above=1pt,hw label]{起始地址} (0.72,-1.5) -- (0.72,-1.86);
  \node[hw label] at (2.65,-3.35) {一次传输中页寄存器不变，缓冲区不得跨越 64 KiB 页边界};
\end{pdfdiagram}
\diagramnote{8237 的 24 位地址拼接：外部页寄存器锁存 A23--A16，通道地址寄存器提供 A15--A0。示例页号 \code{0x02}、页内偏移 \code{0x3450} 合成物理地址 \code{0x23450}；8192 字节传到 \code{0x2544F} 为止，仍在页 \code{0x02}（\code{0x20000}--\code{0x2FFFF}）之内。一次传输中页寄存器保持不变，跨 64 KiB 页边界的缓冲区必须拆成两次传输。}

\begin{lstlisting}
; 8237 ch2 block transfer: device -> memory, 8192 B at phys 0x23450
mov al, 0x06    ; mask ch2 while programming
out 0x0A, al    ; single-channel mask register
mov al, 0x86    ; mode: block, write-to-memory, ch2
out 0x0B, al    ; mode register
out 0x0C, al    ; clear byte-pointer flip-flop (value ignored)
mov al, 0x50    ; address low byte  (0x3450)
out 0x04, al
mov al, 0x34    ; address high byte
out 0x04, al
mov al, 0xFF    ; count low byte    (0x1FFF = 8192-1)
out 0x05, al
mov al, 0x1F    ; count high byte
out 0x05, al
mov al, 0x02    ; page register -> A16..A23
out 0x81, al
mov al, 0x02    ; unmask ch2
out 0x0A, al
; device raises DREQ; 8237 requests the bus and moves bytes;
; TC (count wraps through 0) ends the transfer and raises EOP/IRQ
\end{lstlisting}

序列里有两处容易出错。其一，16 位寄存器经同一个 8 位端口分两次写入，内部字节指针触发器决定下一次写的是低字节还是高字节，配置前必须先发清除命令，否则高低字节会写反。其二，屏蔽必须在配置之前、开放必须在配置之后，否则设备可能在地址只写了一半时就开始搬运。周期窃取的比例可以估算：软盘数据率约 62.5 KB/s，即每 $16\,\mu\text{s}$ 来一个字节；设总线周期 $250\,\text{ns}$、一次单字搬运占 4 拍共 $1\,\mu\text{s}$，DMA 占用总线的时间比例约 $1/16\approx6\%$，其余 94\% 的总线时间仍归 CPU。上面的配置序列按块模式给出：8237 一旦拿到总线就传满 8192 字节，按一次搬运 $1\,\mu\text{s}$ 算是 $8192\times1\,\mu\text{s}\approx8.2\,\text{ms}$ 内总线全部归 DMA。但这个数字只有在设备能持续供数时才成立——块模式传输的跨度实际由 DREQ 决定，软盘每 $16\,\mu\text{s}$ 才来一个字节，传完 8192 字节仍要 $8192\times16\,\mu\text{s}\approx131\,\text{ms}$，并不比单字模式快，而这 131 ms 里总线被 DMA 占着，CPU 长时间取不到指令。正因如此，PC 的软盘 DMA 实际配成单字模式（模式字节 \code{0x46}）；块模式适合的是设备端能跟上总线速度的场合，如内存到内存传送或磁盘内部缓冲命中后的连续读出。选哪种模式，本质上是在搬运总时间与 CPU 停顿之间做权衡。

\topic{DMA 最容易错在可见性和所有权}
DMA 让 CPU 和设备共享主存，也引入 cache 一致性、顺序和缓冲区所有权问题。CPU 写好的描述符必须先对设备可见；设备写入的缓冲区必须在 CPU 读取前变得可见；CPU 不能在设备仍拥有缓冲区时重用它；设备不能越过分配的地址范围。现代系统还会用 IOMMU 限制 DMA 能访问哪些物理页，避免坏设备或错误驱动破坏任意内存。

\begin{pdfdiagram}
  \node[hw state, minimum width=1.85cm] (cpu) at (0,0) {CPU owns};
  \node[hw compute, minimum width=2.1cm] (desc) at (2.65,0) {描述符可见};
  \node[hw control, minimum width=1.95cm] (dev) at (5.35,0) {Device owns};
  \node[hw compute, minimum width=2.0cm] (comp) at (7.95,0) {completion};
  \node[hw state, minimum width=1.85cm] (cpu2) at (10.45,0) {CPU owns};
  \draw[hw bus] (cpu) -- node[above,hw label]{提交描述符} (desc);
  \draw[hw bus] (desc) -- node[above,hw label]{设备搬运} (dev);
  \draw[hw bus] (dev) -- node[above,hw label]{写状态/中断} (comp);
  \draw[hw return] (comp) -- node[above,hw label]{同步 cache 后读取} (cpu2);
  \node[hw label, text width=10.0cm] at (5.25,-1.0) {所有权不能重叠：CPU 交出缓冲区后不得改写；设备 completion 到达且数据可见后，CPU 才能重新使用。};
\end{pdfdiagram}
\diagramnote{DMA 缓冲区所有权。完成通知、cache 可见性和所有权回收必须按顺序发生。}

\begin{longtable}{L{0.18\linewidth}L{0.38\linewidth}L{0.34\linewidth}}
\toprule
边界 & 必须说明 & 失败后果 \\
\midrule
描述符可见 & CPU 写描述符后如何让设备看到 & 设备读到旧地址、旧长度或旧标志 \\
数据可见 & 设备写缓冲区后 CPU 如何看到新数据 & CPU 读到 cache 中的旧内容 \\
所有权 & 缓冲区何时归 CPU，何时归设备 & 重用正在 DMA 的缓冲区，数据被覆盖 \\
地址权限 & 设备能访问哪些物理页或 I/O 虚拟地址 & DMA 写坏无关内存或触发 IOMMU fault \\
完成顺序 & completion 到达是否代表所有数据已可见 & 收到中断后仍读到未完成数据 \\
\bottomrule
\end{longtable}

\section{五、从教学总线到现代互连}

早期微机的并行总线比较直观：地址线、数据线、读写控制、片选和等待信号都能在逻辑分析仪上看到。现代芯片内部和 PCIe 这类外部互连更复杂，可能使用分层 packet、请求队列、completion、错误报告和流控。但抽象问题没有改变：事务由发起者提出，经互连路由到目标，目标返回数据或完成状态，错误路径必须被记录和上报。

\begin{longtable}{L{0.18\linewidth}L{0.36\linewidth}L{0.36\linewidth}}
\toprule
互连对象 & 旧总线中的类比 & 现代形式 \\
\midrule
地址阶段 & 地址线和片选 & 请求 packet、BAR、路由 ID、地址译码 \\
数据阶段 & 数据线和字节使能 & payload、byte enable、burst 或 cache line \\
等待 & READY/WAIT & credit、ready/valid、队列背压 \\
完成 & 采样数据或写周期结束 & completion packet、状态位、中断或 MSI \\
错误 & 总线错误或超时 & poison、fault、AER、IOMMU fault、设备 error \\
\bottomrule
\end{longtable}

\topic{总线仲裁：多个主设备谁先用总线}
共享总线上可以有多个能发起事务的主设备：CPU、DMA 控制器、显示控制器都可能同时要占用地址线和数据线。同一时刻只能有一个主设备驱动总线，因此必须有一种决定“这一轮给谁”的机制，这就是总线仲裁（bus arbitration）。仲裁要同时回答三个问题：裁决花多少时间，谁拥有优先权，低优先级请求会不会永远等不到。经典做法有三种，它们在线数、延迟和公平性之间做出不同取舍。

集中式仲裁器为每个主设备拉一对独立的请求线（request）和授权线（grant），全部连到一个中央仲裁器。所有请求并行到达，仲裁器在一两个周期内按优先级逻辑选出获胜者并抬起对应的 grant。这是最快的做法，优先级还可以做成可编程寄存器；代价是线数随主设备数线性增长，$N$ 个主设备需要 $2N$ 根专用线，仲裁器本身也成为单点。

菊花链（daisy chain）只拉一条授权线，从仲裁器出发依次穿过所有主设备。设备没有请求时把授权信号向下一级传递，有请求时把授权截留下来占用总线。线数与设备数量无关，结构最简单；但优先级由设备在链上的物理位置写死——离仲裁器越近永远越优先，链尾设备可能长期等不到授权，而且授权要逐级传播，链越长延迟越大。

轮询计数（polling）用一小组设备号线代替一对一连线：仲裁器驱动 $\lceil\log_2 N\rceil$ 根线，轮流扫描设备编号，被扫到且有请求的设备抬起 busy 占用总线。线数最少，公平性也最灵活——计数器若每次从零开始，编号小的设备优先；若从上一次授权的设备之后继续扫描，就得到轮转公平。代价是平均要扫过半个编号表才能命中请求者，是三者中最慢的。

\begin{longtable}{L{0.13\linewidth}L{0.15\linewidth}L{0.26\linewidth}L{0.15\linewidth}L{0.21\linewidth}}
\toprule
做法 & 专用线数 & 优先级 & 授权速度 & 关键问题 \\
\midrule
集中式仲裁器 & $2N$（每设备一对请求/授权） & 由仲裁器逻辑决定，可编程 & 最快，并行比较 & 线多，仲裁器是单点 \\
菊花链 & 与设备数无关（一条授权链） & 由物理位置固定，链首最高 & 逐级传播，链长则慢 & 链尾可能等不到，优先级不可改 \\
轮询计数 & $\lceil\log_2 N\rceil$ 根扫描线 & 由计数起点决定，可轮转公平 & 最慢，平均扫半表 & 扫描本身占用总线时间 \\
\bottomrule
\end{longtable}

现代片上互连并没有消灭仲裁，而是把它从“整条共享线上的全局问题”拆小，做进了每一个交换节点。NoC 路由器的每个输入端口要仲裁来自不同方向的 flit，crossbar 要为每个输出端口在多个输入之间仲裁，PCIe switch 要在端口和虚通道之间仲裁。共享并行总线退化成点到点链路之后，“谁先用”的问题依然存在于每个交换机内部——只是裁决范围从整块板卡缩小到一个端口。

\topic{valid/ready 是理解现代片上互连的最小模型}
许多片上接口都可以抽象成 valid/ready 握手。发送方在 valid 为 1 时保持请求或响应稳定，接收方在 ready 为 1 时表示可以接收，只有二者同周期为 1，事务片段才真正完成。请求通道和响应通道可以分开，读请求和读数据也可以隔很多周期。这个模型比“总线一拍完成”更接近现代硬件。

\begin{pdfdiagram}[x=1cm,y=1cm]
  % valid — 先升高并等待
  \draw[BookTeal,line width=0.85pt] (0.6,0.45) -- (1.0,0.45) -- (1.0,1.2) -- (5.5,1.2) -- (5.5,0.45) -- (9.0,0.45);
  \node[hw label, anchor=east] at (0.5,1.2) {valid};
  % ready — 背压阶段为 0, 准备好后才升高
  \draw[BookAccent,line width=0.85pt] (0.6,0.0) -- (4.0,0.0) -- (4.0,0.75) -- (5.5,0.75) -- (5.5,0.0) -- (9.0,0.0);
  \node[hw label, anchor=east] at (0.5,0.38) {ready};
  % payload — valid=1 期间保持稳定的总线
  \draw[BookMuted,line width=0.65pt,dashed] (0.6,-0.75) -- (1.0,-0.75);
  \draw[BookMuted,line width=0.65pt,dashed] (5.5,-0.75) -- (9.0,-0.75);
  \draw[BookTeal,line width=0.85pt] (1.0,-0.45) -- (5.5,-0.45) -- (5.5,-1.05) -- (1.0,-1.05) -- cycle;
  \node[hw label, anchor=east] at (0.5,-0.75) {payload};
  \node[hw label] at (3.25,-0.75) {数据保持稳定};
  \node[hw label] at (7.25,-0.75) {允许变化};
  % 区间竖线与 transfer
  \draw[BookRule,line width=0.45pt,dashed] (1.0,1.6) -- (1.0,-1.3);
  \draw[BookRule,line width=0.45pt,dashed] (4.0,1.6) -- (4.0,-1.3);
  \draw[BookRule,line width=0.45pt,dashed] (5.5,1.6) -- (5.5,-1.3);
  \node[hw label] at (2.5,1.85) {背压：valid 等待};
  \node[hw label] at (4.75,1.85) {transfer};
  \node[hw label, text width=8.4cm] at (4.7,-1.75) {valid=1 且 ready=0 时，发送方必须保持 payload 稳定；二者同周期为 1 才完成一次传输。};
\end{pdfdiagram}
\diagramnote{valid/ready 握手。发送方先抬 valid 并保持数据稳定；接收方在背压期间保持 ready=0；两条线在 transfer 周期同时为 1，事务片段才完成。背压不是错误，而是接收方暂时不能接受。}

\begin{lstlisting}
handshake rule:
  transfer happens only when valid == 1 and ready == 1
  while valid == 1 and ready == 0:
    sender keeps payload stable
  receiver may apply backpressure by deasserting ready
\end{lstlisting}

\topic{三种 valid/ready 时序情形}
valid/ready 的全部行为可以归纳成一条采样规则和三种常见情形。采样规则是：只在某个采样沿上 valid 和 ready 同时为 1，该通道才完成一次传输。三种常见情形是：双方同拍就绪，当拍完成；接收方拉低 ready 施加背压，发送方必须保持 valid 和 payload 不变；请求通道和响应通道各自独立握手，两次传输之间可以隔任意多个周期。这三种情形都不是异常，而是流控协议的正常组成部分。

\begin{waveform}[x=0.8cm,y=0.8cm]
  % (a) 同拍传输
  \node[hw label, anchor=west] at (0,11.3) {(a) 同拍传输};
  \draw[BookTeal, line width=0.9pt] (0,10.2) -- (4,10.2) -- (4,10.9) -- (8,10.9) -- (8,10.2) -- (12,10.2);
  \node[hw label, anchor=east] at (-0.1,10.55) {valid};
  \draw[BookAccent, line width=0.9pt] (0,9.2) -- (4,9.2) -- (4,9.9) -- (8,9.9) -- (8,9.2) -- (12,9.2);
  \node[hw label, anchor=east] at (-0.1,9.55) {ready};
  \draw[BookMuted, line width=0.65pt, dashed] (0,8.55) -- (4,8.55);
  \draw[BookMuted, line width=0.65pt, dashed] (8,8.55) -- (12,8.55);
  \draw[BookTeal, line width=0.9pt] (4,8.2) rectangle (8,8.9);
  \node at (6,8.55) {D0};
  \node[hw label, anchor=east] at (-0.1,8.55) {payload};
  \draw[BookRule, line width=0.45pt, dashed] (8,11.2) -- (8,8.05);
  \node[hw label, anchor=south west] at (8.1,11.25) {传输沿};
  % (b) 背压保持
  \node[hw label, anchor=west] at (0,7.55) {(b) 背压保持};
  \draw[BookTeal, line width=0.9pt] (0,6.45) -- (4,6.45) -- (4,7.15) -- (12,7.15) -- (12,6.45);
  \node[hw label, anchor=east] at (-0.1,6.8) {valid};
  \draw[BookAccent, line width=0.9pt] (0,5.45) -- (8,5.45) -- (8,6.15) -- (12,6.15) -- (12,5.45);
  \node[hw label, anchor=east] at (-0.1,5.8) {ready};
  \draw[BookMuted, line width=0.65pt, dashed] (0,4.8) -- (4,4.8);
  \draw[BookTeal, line width=0.9pt] (4,4.45) rectangle (12,5.15);
  \node at (8,4.8) {D1（保持不变）};
  \node[hw label, anchor=east] at (-0.1,4.8) {payload};
  \draw[BookRule, line width=0.45pt, dashed] (4,7.45) -- (4,4.3);
  \node[hw label, anchor=south west] at (4.1,7.5) {valid 抬起};
  \draw[BookRule, line width=0.45pt, dashed] (12,7.45) -- (12,4.3);
  \node[hw label, anchor=south east] at (11.9,7.5) {传输沿};
  % (c) 请求响应分离
  \node[hw label, anchor=west] at (0,3.75) {(c) 请求响应分离};
  \draw[BookTeal, line width=0.9pt] (0,3.35) -- (4,3.35) -- (4,2.65) -- (12,2.65);
  \node[hw label, anchor=east] at (-0.1,3.0) {req valid};
  \draw[BookTeal, line width=0.9pt] (0,1.65) rectangle (4,2.35);
  \node at (2,2.0) {请求};
  \draw[BookMuted, line width=0.65pt, dashed] (4,2.0) -- (12,2.0);
  \node[hw label, anchor=east] at (-0.1,2.0) {req data};
  \draw[BookTeal, line width=0.9pt] (0,0.65) -- (8,0.65) -- (8,1.35) -- (12,1.35) -- (12,0.65);
  \node[hw label, anchor=east] at (-0.1,1.0) {resp valid};
  \draw[BookAccent, line width=0.9pt] (0,0.35) -- (12,0.35);
  \node[hw label, anchor=east] at (-0.1,0.0) {resp ready};
  \draw[BookMuted, line width=0.65pt, dashed] (0,-1.0) -- (8,-1.0);
  \draw[BookTeal, line width=0.9pt] (8,-1.35) rectangle (12,-0.65);
  \node at (10,-1.0) {响应};
  \node[hw label, anchor=east] at (-0.1,-1.0) {resp data};
  \draw[BookRule, line width=0.45pt, dashed] (4,3.65) -- (4,1.5);
  \node[hw label, anchor=south west] at (4.1,3.7) {请求传输};
  \draw[BookRule, line width=0.45pt, dashed] (12,1.65) -- (12,-1.5);
  \node[hw label, anchor=south east] at (11.9,1.7) {响应传输};
  \draw[BookMuted, line width=0.6pt, <->] (4.3,2.95) -- (7.7,2.95);
  \node[hw label, above] at (6,2.95) {隔若干周期};
\end{waveform}
\diagramnote{三种情形一条规则：只在 valid 与 ready 同拍为 1 的沿上传输。(a) 双方同拍就绪，当拍完成；(b) ready 为 0 的期间 valid 与 D1 保持不变，这是接收方施加的背压，不是错误；(c) 请求在一个沿完成，响应通道隔若干周期后才完成自己的传输，两个通道互不等待。}

\topic{写缓冲和内存顺序解释了现代芯片为什么看起来更快}
写入比读取更容易被隐藏。CPU 可以先把写事务放进写缓冲，让执行流水线继续推进；写缓冲再把相邻写合并、排队发送到 cache 或互连。这样常见 STORE 的表面延迟很低，但代价是可见性变复杂：其他核心、DMA 或设备什么时候看到这次写，取决于 cache 一致性、内存类型和屏障规则。普通 RAM、MMIO 和 DMA 描述符不能使用同一套可见性假设。

\topic{总线错误和超时必须成为规格的一部分}
未映射地址不能只写成“不要访问”。真实系统必须规定访问未映射区域时会发生什么：返回固定值、保持等待直到超时、产生总线错误，或进入异常入口。设备无响应、DRAM 控制器报错、ECC 检测失败、外设权限不允许，也都应该形成可观察错误状态。没有错误约定的整机，调试时会把硬件故障误判为程序死循环。

\topic{cache 行和局部性把常见访问留在近处}
主存容量大但离 CPU 远，寄存器离 CPU 近但容量小，cache 处在二者之间。cache 不是“更快的内存副本”这么简单，而是按 cache line 搬运的一组近处 SRAM。CPU 读取某个地址时，cache 会用地址中的 tag 判断目标行是否已经在近处；若命中，数据直接从 cache 返回；若未命中，cache 控制器向下一级 cache 或 DRAM 请求整行，再把需要的字节送回 CPU。

\begin{longtable}{L{0.18\linewidth}L{0.36\linewidth}L{0.36\linewidth}}
\toprule
对象 & 硬件含义 & 关键问题 \\
\midrule
cache line & 一次填充和替换的最小数据块 & 行大小、字节偏移和未对齐访问是否清楚 \\
tag & 判断某个 cache 槽保存的是哪段地址 & 命中比较是否覆盖高位地址和有效位 \\
index & 选择 cache 中的候选槽或集合 & 映射冲突是否可能造成反复 miss \\
valid/dirty & 表示行是否有效、是否被修改过 & 替换时是否需要写回下一级 \\
miss & 近处没有所需行，需要向下一级请求 & miss 期间 CPU 停顿、继续执行或乱序等待的边界 \\
\bottomrule
\end{longtable}

局部性解释了 cache 为什么有效。时间局部性表示刚访问过的数据很可能再次访问，例如循环变量和栈顶对象；空间局部性表示访问某个地址后，很可能访问相邻地址，例如顺序扫描数组。cache line 正是利用空间局部性：一次 miss 不只取一个字节，而是取一整段相邻字节。代价也随之出现：若程序跨很大步长访问数组，或者两个热数据反复映射到同一 cache 集合，硬件仍会不停搬行，表面上只是 LOAD 很慢，实质上是事务层反复 miss 和替换。

\begin{lstlisting}
32-byte cache line example:
  address bits = [ tag | index | byte_offset ]
  byte_offset selects one byte within the line
  index selects a cache set
  tag comparison decides hit or miss
\end{lstlisting}

\topic{案例：一次 cache 命中的完整分解}cache 的查找可以拆成三个字段各管一件事：offset 选行内字节，index 选槽，tag 判身份。以 32-bit 地址、16 KiB 直接映射、64 B 行为例：offset $=\log_2 64=6$ bit；行数 $=16\ \text{KiB}\div 64\ \text{B}=256$，index $=\log_2 256=8$ bit；tag $=32-8-6=18$ bit。任何地址到来时，硬件先按 index 找到唯一的槽，再用 tag 和 valid 位判断“槽里这行是不是这段地址”。

\begin{pdfdiagram}
  % 32-bit 地址三段
  \draw[hw bus] (0,4.0) rectangle (4.0,4.6);
  \node at (2.0,4.3) {tag（18 bit）};
  \draw[hw bus] (4.0,4.0) rectangle (6.4,4.6);
  \node at (5.2,4.3) {index（8 bit）};
  \draw[hw bus] (6.4,4.0) rectangle (8.6,4.6);
  \node at (7.5,4.3) {offset（6 bit）};
  \node[hw label, above] at (4.3,4.6) {32-bit 地址（16 KiB 直接映射、64 B 行）};
  % tag → 比较器
  \draw[hw bus] (2.0,4.0) -- (2.0,2.32);
  \node[hw compute, minimum width=2.4cm] (cmp) at (1.5,2.05) {tag 比较（= ?）};
  % index → 槽
  \draw[hw bus] (5.6,4.0) -- node[hw label, right, pos=0.55]{index 译码选槽} (5.6,3.18);
  \node[hw state, minimum width=4.8cm] (s0) at (5.6,2.9) {槽 168：tag，V，数据};
  \node[hw state, minimum width=4.8cm] (s1) at (5.6,2.05) {槽 169：tag=0x00010，V=1};
  \node[hw state, minimum width=4.8cm] (s2) at (5.6,1.2) {槽 170：tag，V，数据};
  \draw[hw bus] (s1.west) -- (cmp.east);
  \draw[hw bus] (cmp.south) -- (1.5,0.7) node[hw label, right, pos=1.0, text width=3.6cm]{hit：读槽内数据；miss：向下一级取整行替换};
  % offset → 字节选择
  \draw[hw bus] (7.5,4.0) -- (7.5,3.3) node[hw label, right, pos=0.6, text width=3.4cm]{offset 选行内字节（0x15 → 第 21 字节）};
\end{pdfdiagram}
\diagramnote{直接映射 cache 的查找。地址拆成 tag/index/offset 三段：index 译码选出唯一的槽，槽内 tag 与地址 tag 比较决定 hit/miss，offset 再从 64 B 行内选出目标字节。}

下表把四次访问的分解和结果走一遍。注意第 3、4 次展示了直接映射最典型的失效模式：

\begin{longtable}{L{0.08\linewidth}L{0.17\linewidth}L{0.28\linewidth}L{0.39\linewidth}}
\toprule
次序 & 地址 & tag / index / offset & 结果 \\
\midrule
1 & \code{0x00042A40} & \code{0x00010} / \code{0xA9} / \code{0x00} & 槽 169 的 V=0 → miss；从 DRAM 取回 64 B 整行填入，tag 写 \code{0x00010}、V 置 1，返回第 0 字节 \\
2 & \code{0x00042A40} & \code{0x00010} / \code{0xA9} / \code{0x00} & tag 相等且 V=1 → hit，直接从槽内返回数据 \\
3 & \code{0x00082A55} & \code{0x00020} / \code{0xA9} / \code{0x15} & 槽 169 的 tag 是 \code{0x00010}，不等 → 冲突 miss；取回新行替换旧行（旧行若 dirty 先写回 DRAM），返回第 21 字节 \\
4 & \code{0x00042A40} & \code{0x00010} / \code{0xA9} / \code{0x00} & 槽 169 刚被第 3 次访问换掉 → 又 miss \\
\bottomrule
\end{longtable}

两个热地址映射到同一个槽时，直接映射会反复 miss——这正是组相联存在的理由：同一个槽放多路，几路 tag 并行比较，任何一路相等都算 hit，常用数据就不必互相驱逐。

\topic{两路组相联让热行各占一路}
直接映射的冲突来自“一个槽只能放一行”。把同样 16 KiB、64 B 行的 cache 改成 128 组、每组 2 路：offset 仍是 6 bit，index 从 8 bit 变成 7 bit，tag 从 18 bit 变成 19 bit。重走前面的四次访问：\code{0x00042A40} 和 \code{0x00082A55} 的 index 都是 \code{0x29}，tag 分别是 \code{0x00021} 和 \code{0x00041}。第 1 次 miss，填入路 0；第 2 次 tag 相等，hit；第 3 次两路 tag 都不等（路 1 的 V 仍为 0），仍然 miss，但新行填入路 1，不再驱逐路 0；第 4 次再读 \code{0x00042A40}，路 0 的 tag 相等，变成 hit。直接映射下第 3、4 次的互相驱逐，在 2 路组相联下只剩第 3 次那次无法避免的首次 miss。代价是每组多一个 tag 比较器、一个数据选择 mux 和若干替换位。

\begin{pdfdiagram}
  % 32-bit 地址三段（2 路组相联：tag 19 / index 7 / offset 6）
  \draw[hw bus] (0,5.6) rectangle (3.2,6.2);
  \node at (1.6,5.9) {tag（19 bit）};
  \draw[hw bus] (3.2,5.6) rectangle (5.6,6.2);
  \node at (4.4,5.9) {index（7 bit）};
  \draw[hw bus] (5.6,5.6) rectangle (7.8,6.2);
  \node at (6.7,5.9) {offset（6 bit）};
  \node[hw label, above] at (3.9,6.2) {32-bit 地址（16 KiB、2 路组相联、64 B 行）};
  % index 译码：干线加两路下引线
  \draw[hw bus] (4.4,5.6) -- (4.4,4.9);
  \draw[hw bus] (2.0,4.9) -- (8.3,4.9);
  \node[hw label, below] at (5.15,4.9) {index 译码：同一组的两路同时选中};
  \draw[hw bus] (2.0,4.9) -- (2.0,3.91);
  \draw[hw bus] (8.3,4.9) -- (8.3,3.91);
  % offset 选行内字节
  \node[hw label, below] at (6.7,5.6) {offset 选行内字节};
  % 组内两路
  \node[hw state, minimum width=3.4cm] (way0) at (2.0,3.6) {路 0：tag，V，64 B 数据};
  \node[hw state, minimum width=3.4cm] (way1) at (8.3,3.6) {路 1：tag，V，64 B 数据};
  % 2 选 1 mux
  \node[hw compute, minimum width=2.4cm] (mux) at (5.15,3.6) {2 选 1 mux};
  \draw[hw bus] (way0.east) -- (mux.west);
  \draw[hw bus] (way1.west) -- (mux.east);
  \draw[hw bus] (5.55,3.29) -- (5.55,2.85);
  \node[hw label, anchor=east] at (5.05,3.05) {命中路数据};
  % 两路 tag 并行比较
  \node[hw compute, minimum width=2.4cm] (cmp0) at (2.0,2.35) {tag 比较（= ?）};
  \node[hw compute, minimum width=2.4cm] (cmp1) at (8.3,2.35) {tag 比较（= ?）};
  \draw[hw bus] (way0.south) -- (cmp0.north);
  \node[hw label, anchor=west] at (2.05,2.72) {槽内 tag};
  \draw[hw bus] (way1.south) -- (cmp1.north);
  \node[hw label, anchor=west] at (8.35,2.72) {槽内 tag};
  % 地址 tag 总线：并行送两个比较器
  \draw[hw bus] (1.6,5.6) -- (1.6,5.05) -- (0.3,5.05) -- (0.3,0.9) -- (8.0,0.9) -- (cmp1.south);
  \node[hw label, above] at (0.95,5.05) {地址 tag};
  \draw[hw bus] (0.3,2.35) -- (cmp0.west);
  \node[hw label, below] at (3.2,0.9) {地址 tag 并行送两路比较器};
  % 任一路相等即 hit，并选择 mux
  \node[hw compute, minimum width=1.6cm] (orn) at (5.15,2.35) {或（hit）};
  \draw[hw bus] (cmp0.east) -- (orn.west);
  \draw[hw bus] (cmp1.west) -- (orn.east);
  \draw[hw return] (orn.north) -- (mux.south);
  \node[hw label, anchor=west] at (5.35,2.9) {hit / 选路};
  % 替换位：miss 时选一路驱逐
  \node[hw state, minimum width=1.6cm] (rep) at (5.15,1.35) {替换位};
  \draw[hw return] (rep.north) -- (orn.south);
  \node[hw label, anchor=west] at (5.25,1.7) {miss 时选一路驱逐};
\end{pdfdiagram}
\diagramnote{2 路组相联的查找。index 译码选中整个组，组内两路的 tag 同时与地址 tag 比较；任一路相等，或门即输出 hit，并用比较结果控制 mux 选出命中路的数据。miss 时按组内替换位选一路驱逐、填入新行。}

\topic{写策略决定 STORE 什么时候真正离开 CPU}
读 miss 的动作比较直观：向下级取回数据。写入更复杂，因为 CPU 可以先改 cache，再决定什么时候把修改传播到下一级。写穿策略每次写命中都继续写下一级，行为直接但带宽压力大；写回策略只在 cache 中标脏位，替换时再写回下一级，常见情况下更快，但可见性和错误恢复更复杂。无论哪种策略，字节写都必须尊重 byte enable，不能把同一 cache line 中无关字节改坏。

\begin{longtable}{L{0.18\linewidth}L{0.36\linewidth}L{0.36\linewidth}}
\toprule
写策略 & 优点 & 风险和边界 \\
\midrule
写穿 & 下一级更快看见新值，调试直观 & 带宽压力大，常配合写缓冲隐藏延迟 \\
写回 & 多次写同一行可合并，减少外部事务 & dirty 行替换、掉电和 DMA 可见性更难处理 \\
写分配 & 写 miss 时先把整行取入 cache 再修改 & 小写入可能触发整行读取 \\
非写分配 & 写 miss 直接向下级写，不装入 cache & 后续读同地址可能仍 miss \\
\bottomrule
\end{longtable}

这也是 MMIO 不能随意缓存的根本原因。写 LED 寄存器、清中断状态或启动 DMA 命令时，程序需要的是一次外设可见动作，而不是某条 cache line 中的普通字节更新。若 MMIO 被写回 cache 吞掉，设备可能永远看不到命令；若读状态寄存器被 cache 命中，CPU 可能一直看到旧状态。因此地址属性必须区分普通可缓存内存、不可缓存设备内存、写合并 framebuffer、DMA coherent 区域等不同事务约定。

\topic{cache miss 是事务暂停，不是抽象的慢}
当 LOAD miss 时，流水线不能凭空得到数据。简单 CPU 会停在访存阶段，保持地址和控制状态，直到下级返回数据；高性能 CPU 可能让独立指令继续执行，把这条 LOAD 放在未完成队列中，等数据回来再唤醒依赖指令。二者复杂度不同，但都必须维护同一个可见性约定：依赖这次 LOAD 的指令不能使用错误旧值，异常和中断也不能让程序看见半提交状态。

\begin{longtable}{L{0.18\linewidth}L{0.38\linewidth}L{0.34\linewidth}}
\toprule
miss 类型 & 事务路径 & 典型处理 \\
\midrule
L1 miss/L2 hit & L1 请求 L2 返回 cache line & 几到十几周期，流水线可能短暂停顿 \\
最后级 cache miss & 请求内存控制器和 DRAM & 上百周期级别，常需乱序、预取或多线程隐藏 \\
TLB miss & 地址转换缓存缺失，要查页表或进异常 & 页表 walker、权限检查和可能的页错误 \\
写回替换 & 脏行被换出，需先写下一级 & 替换队列、写缓冲和错误报告 \\
设备/DMA 冲突 & cache 中旧副本与设备写入不一致 & coherent 互连、cache 维护或不可缓存映射 \\
\bottomrule
\end{longtable}

\topic{低时延来自近处、缓存、并行和卸载}
现代整机降低时延，不是让所有路径都变成零等待。L1/L2 cache 让常见数据走近处 SRAM；TLB 让近期地址转换不必每次查页表；预取和分支预测提前准备可能需要的路径；流水线、旁路和乱序调度让独立工作重叠；DMA、队列和中断聚合把批量 I/O 从 CPU 指令流中卸载。失败路径仍然存在：cache miss、TLB miss、DRAM 排队、设备背压、DMA fault 和错误中断都会把访问拉回慢路径。

\topic{六、存储层次、地址转换和设备驱动 trace（执行轨迹）}

前面已经讲了单次总线事务、MMIO、中断、DMA 和 cache。真正的系统程序运行时，这些对象不会孤立出现。一次普通 LOAD 可能先经过 TLB 查找，再查 L1 cache，未命中后查 L2/L3，最后进入内存控制器；一次设备接收数据可能由驱动准备描述符，设备 DMA 写缓冲区，cache 一致性或维护操作让 CPU 看见新数据，中断再把完成事件带回来。若教材只讲“内存比寄存器慢”，就漏掉了真正的硬件路径。

\begin{longtable}{L{0.18\linewidth}L{0.4\linewidth}L{0.32\linewidth}}
\toprule
访问层次 & 常见命中路径 & 失败或慢路径 \\
\midrule
寄存器 & 直接从寄存器堆读写 & 端口冲突、旁路失败或流水线停顿 \\
L1 cache & 近处 SRAM 命中，几个周期内返回 & L1 miss，向 L2 请求整条 cache line \\
L2/L3 cache & 片上较大 cache 命中 & 最后级 miss，访问内存控制器和 DRAM \\
TLB & 地址转换缓存命中 & page walk、页错误或权限异常 \\
DRAM & 控制器调度行命中和突发传输 & 行冲突、刷新、排队、ECC 错误 \\
设备/MMIO & 设备寄存器响应读写 & 背压、超时、错误状态或中断完成 \\
\bottomrule
\end{longtable}

\topic{TLB 把地址转换也纳入近处缓存}
虚拟地址或线性地址到物理地址的转换若每次都查页表，访存会变得很慢。TLB 可以理解为“地址转换的 cache”：它保存近期虚拟页到物理页框的映射及权限位。LOAD/STORE 形成虚拟地址后，先用虚拟页号查 TLB；命中时得到物理页框和权限信息，再与页内偏移组合成物理地址；未命中时，硬件 page walker 或软件异常处理会去内存里的页表结构查找。

\begin{lstlisting}
virtual address = [ virtual_page_number | page_offset ]
TLB hit:
  physical_address = TLB[virtual_page_number].frame | page_offset
TLB miss:
  walk page tables or enter miss handler
\end{lstlisting}

\begin{longtable}{L{0.18\linewidth}L{0.36\linewidth}L{0.36\linewidth}}
\toprule
TLB 字段 & 作用 & 关键问题 \\
\midrule
虚拟页号 tag & 判断当前地址是否命中某个转换项 & 进程切换、地址空间标识和刷新规则是否清楚 \\
物理页框 & 与页内偏移组合成物理地址 & 页大小、对齐和物理地址宽度是否一致 \\
权限位 & 读写、执行、用户/特权等访问属性 & 访问失败是权限异常，不是普通 cache miss \\
valid/global & 表示转换是否有效、是否跨进程保留 & 页表修改后如何让旧 TLB 项失效 \\
\bottomrule
\end{longtable}

TLB miss 和 cache miss 都是“近处没有”，但语义不同。cache miss 只是数据不在近处，通常可以向下级取回；TLB miss 可能只是转换缓存缺失，也可能说明页不存在或权限不允许。若页表项不存在，硬件不能继续访问内存数据；若权限不允许，访问必须转入异常。这个差别解释了为什么地址转换、保护和异常在系统硬件中是同一条路径的不同分支。

\topic{多核和 DMA 让 cache 一致性成为硬件约定}
单核教学机里，cache 可以先理解为 CPU 与主存之间的近处副本。多核或 DMA 加入后，问题变成“多个观察者如何看见同一地址的值”。一个核心的 cache 中可能有某条 line 的旧副本，另一个核心或设备刚刚写入了新值；如果没有一致性协议、cache 维护或不可缓存映射，CPU 读到的就可能不是系统真实最新值。

\begin{longtable}{L{0.2\linewidth}L{0.36\linewidth}L{0.34\linewidth}}
\toprule
场景 & 风险 & 常见硬件/软件边界 \\
\midrule
核心 A 写，核心 B 读 & B 的 cache 中可能有旧行 & cache coherence、失效消息、共享/独占状态 \\
CPU 写描述符，设备读 & 描述符可能还在 CPU cache 或写缓冲中 & coherent DMA、flush、写屏障 \\
设备写缓冲区，CPU 读 & CPU cache 中可能有旧数据 & invalidate、coherent 互连、不可缓存映射 \\
MMIO 状态寄存器 & cache 命中会隐藏设备新状态 & 设备内存属性、禁止普通缓存 \\
\bottomrule
\end{longtable}

一致性协议细节可以很深，但教学阶段先抓住三件事。第一，同一地址可能在多个近处缓存中有副本。第二，写入必须让其他观察者的旧副本失效或更新。第三，设备并不总是自动参与 CPU cache 一致性；若平台不支持 coherent DMA，驱动必须在交出缓冲区前后做 cache 维护。否则 DMA 完成中断到达后，CPU 仍可能读到旧 cache line。

\topic{MESI 用四个状态约定谁拿着最新值}
这套状态存在的意义，是让任何时刻、任何地址在所有观察者眼里只有一份“最新值”。想写某一行，核必须先成为唯一持有者：从 S 升级要广播失效，把其他核的副本打成 I；只有处于 M 或 E 时才能直接改。想读可以共享：多个核同时持有 S，读到的都是同一个干净副本，但共享期间谁也不能直接写。于是两个核不会各自拿着不同的新值——要么一个核独占、其余无效，要么大家一起只读。参与一致性互连的 DMA 设备也被同一套状态约束；不参与的平台，驱动就要用 cache clean/invalidate 手工扮演“写回”和“失效”的角色。

\begin{pdfdiagram}
  \node[hw state, minimum width=1.9cm] (stI) at (1.6,3.2) {I 无效};
  \node[hw state, minimum width=1.9cm] (stS) at (9.6,3.2) {S 共享};
  \node[hw state, minimum width=1.9cm] (stE) at (1.6,0.5) {E 独占};
  \node[hw state, minimum width=1.9cm] (stM) at (9.6,0.5) {M 已修改};
  % 本核读：I 分裂为 S 或 E
  \draw[hw return, ->] (stI.east) -- node[above, hw label] {本核读：发现他人副本} (stS.west);
  \draw[hw return, ->] (stI.south) -- (stE.north);
  \node[hw label, anchor=west] at (1.7,1.85) {本核读：无其他副本};
  % 本核写与侦听他人读：S 与 M 之间
  \draw[hw return, <->] (stS.south) -- (stM.north);
  \node[hw label, anchor=east] at (9.5,2.2) {本核写：广播失效 $\rightarrow$ M};
  \node[hw label, anchor=east] at (9.5,1.5) {他人读：写回 $\rightarrow$ S};
  % 本核写：E 已独占，直接改
  \draw[hw return, ->] (stE.east) -- node[above, hw label] {本核写：已独占，直接改} (stM.west);
  % 侦听他人写：副本失效
  \draw[hw return, ->] (stS.north) -- (9.6,3.95) -- (1.6,3.95) -- (stI.north);
  \node[hw label, above] at (5.6,3.95) {侦听他人写：副本失效};
\end{pdfdiagram}
\diagramnote{MESI 一致性状态图（教学简化版）。本核读把 I 变成 E（无其他副本）或 S（他人有副本）；本核写必须先独占：S 到 M 要广播失效，E 到 M 直接改；总线侦听到他人读时 M 写回并退到 S；侦听到他人写时副本失效（E、M 同样失效，M 先写回）。}

\topic{案例：MESI 在两个核之间的完整走查}
状态图只有放进具体访问序列才有意义。下面用两个核、一个地址 X 把上一专题的状态机完整走一遍。设初始时核 A、核 B 的 cache 中 X 行均为 I（无效），内存持有 X 的最新值。先说明一个细节：按严格 MESI，核 A 第一次读入 X 时若没有发现任何其他副本，应记为 E（独占）；本走查假设 X 此前已被第三个核读过，或采用“读入即共享”的教学约定，因此第一步记为 S。这不影响后续步骤的转换逻辑。

\begin{longtable}{L{0.06\linewidth}L{0.24\linewidth}L{0.09\linewidth}L{0.09\linewidth}L{0.42\linewidth}}
\toprule
步骤 & 动作 & 核 A & 核 B & 互连上发生的事 \\
\midrule
0 & 初始 & I & I & X 的最新值在内存中 \\
1 & 核 A 读 X & S & I & A 的读 miss，经互连从内存取回整行，记为共享副本 \\
2 & 核 B 读 X & S & S & B 的读 miss，从内存取行；A 侦听到这次读，保持 S，两核持有同一份干净副本 \\
3 & 核 A 写 X & M & I & A 先广播失效消息，B 的副本作废（S $\rightarrow$ I）；确认后 A 独占修改本行（S $\rightarrow$ M），内存中的 X 从此过期 \\
4 & 核 B 读 X & S & S & B 的读 miss；A 侦听到读，把脏行写回内存（或直接转发给 B），自己降为 S；B 得到最新数据的共享副本 \\
\bottomrule
\end{longtable}

逐步看两处关键转换。第 3 步是“写之前先独占”：A 不能只改自己手里的 S 副本，否则 B 会继续读到旧值，同一地址就出现了两个“最新值”；广播失效把其他副本全部打成 I 之后，A 才升级成 M 并修改。第 4 步是“最新值跟着 M 走”：此时内存里的 X 是旧数据，B 的读不能由内存直接回答，必须由持有 M 行的 A 先把数据交出来——写回内存或 cache 之间直接转发——之后两核各持 S，内存恢复最新。

这张表就是一致性协议的 trace：每一行对应状态图上的一条边，每一步结束都满足同一条不变式——X 的最新值在全局唯一，要么在某核的 M 行里（其余全为 I，内存过期），要么在所有 S 副本与内存里保持一致。调试多核程序时若怀疑数据不一致，最有效的办法正是把可疑的访问序列按这种方式列成表，逐步检查哪一步破坏了唯一性。

\topic{store buffer 和屏障把“执行完成”与“可见”分开}
STORE 对流水线来说可以很早“执行完成”：地址和数据已经形成，写事务进入 store buffer，后续独立指令可以继续执行。但这个 STORE 什么时候对其他核心、DMA 或设备可见，要看 store buffer 排队、cache 状态、内存类型和屏障规则。多核内存模型本身超出本书范围，这里只需要与 DMA 直接相关的一条规则：普通 RAM 写可以合并和延迟；MMIO 写通常要求更强顺序；DMA 描述符写在通知设备前必须已经对设备可见。

\begin{lstlisting}
prepare DMA descriptor:
  descriptor.addr = buffer
  descriptor.len  = length
  memory_barrier()
  MMIO[DEVICE_DOORBELL] = descriptor_index
\end{lstlisting}

这里的屏障不是“让 CPU 慢一点”的软件仪式，而是要求硬件在门铃 MMIO 写对设备可见之前，先让描述符内容按平台规则可见。若顺序反了，设备可能看到门铃，却读到旧描述符或半写入描述符。许多难复现的设备错误，本质上就是这个可见性边界没有写清楚。

\begin{longtable}{L{0.18\linewidth}L{0.38\linewidth}L{0.34\linewidth}}
\toprule
对象 & 可见性要求 & 失败表现 \\
\midrule
普通数据 & 由语言、编译器和内存模型共同约束 & 多核读到旧值或乱序观察 \\
锁变量 & 获取/释放顺序必须约束临界区数据 & 进入临界区后仍看见旧数据 \\
DMA 描述符 & 通知设备前必须对设备可见 & 设备搬错地址或长度 \\
MMIO 门铃 & 前面的配置写必须先到设备接口 & 设备按旧配置启动 \\
中断完成 & completion 到达后数据必须已可读取或可同步 & 中断处理读到旧缓冲区 \\
\bottomrule
\end{longtable}

\topic{案例：一个内存屏障挡住什么}
屏障的作用可以用一个经典的双核序列看清。两个核共享变量 x 和 y，初值都为 0：核 A 先把 1 写入 x，再读 y 到 r1；核 B 先把 1 写入 y，再读 x 到 r2。直觉上，两次写和两次读无论怎样交错，总有一个读应该看到对方的新值 1。但在带 store buffer 的真实机器上，r1=0 且 r2=0 的结果确实可能出现。

\begin{lstlisting}
core A                    core B
mov qword ptr [x], 1      mov qword ptr [y], 1
mov rax, [y]              mov rbx, [x]
\end{lstlisting}

原因就在上一专题的可见性划分：\code{x=1} 对核 A 来说“执行完成”只表示写事务进入了 A 的 store buffer，并不意味着其他核已经能看见它。两个核的 store 都可能还排在各自的 store buffer 里，两个 load 就已经从内存（或对方尚未更新的 cache 副本）里取回了 0。

\begin{longtable}{L{0.08\linewidth}L{0.30\linewidth}L{0.30\linewidth}L{0.22\linewidth}}
\toprule
次序 & 核 A & 核 B & 内存中可见的 x / y \\
\midrule
1 & x=1 进入 A 的 store buffer & --- & 0 / 0 \\
2 & --- & y=1 进入 B 的 store buffer & 0 / 0 \\
3 & r1=y，读到 0 & --- & 0 / 0 \\
4 & --- & r2=x，读到 0 & 0 / 0 \\
5 & store buffer 排空，x=1 全局可见 & store buffer 排空，y=1 全局可见 & 1 / 1 \\
\bottomrule
\end{longtable}

把上表的关键时刻画成四张快照，写滞留与读旧值各处于什么位置就更直观。

\begin{pdfdiagram}
  % 表头
  \node[hw label] at (1.6,0.62) {store buffer A};
  \node[hw label] at (4.0,0.62) {内存};
  \node[hw label] at (6.4,0.62) {store buffer B};
  % ① 初始
  \node[hw label, anchor=east] at (-0.65,0) {① 初始};
  \node[hw state, minimum width=0.95cm] (a1) at (0,0) {核 A};
  \node[hw control, minimum width=1.0cm] (ba1) at (1.6,0) {（空）};
  \node[hw memory, minimum width=2.0cm] (m1) at (4.0,0) {x = 0，y = 0};
  \node[hw control, minimum width=1.0cm] (bb1) at (6.4,0) {（空）};
  \node[hw state, minimum width=0.95cm] (b1) at (8.0,0) {核 B};
  % ② 写滞留
  \node[hw label, anchor=east] at (-0.65,-1.5) {② 写滞留};
  \node[hw state, minimum width=0.95cm] (a2) at (0,-1.5) {核 A};
  \node[hw control, minimum width=1.0cm] (ba2) at (1.6,-1.5) {x = 1};
  \node[hw memory, minimum width=2.0cm] (m2) at (4.0,-1.5) {x = 0，y = 0};
  \node[hw control, minimum width=1.0cm] (bb2) at (6.4,-1.5) {y = 1};
  \node[hw state, minimum width=0.95cm] (b2) at (8.0,-1.5) {核 B};
  \draw[hw bus] (a2) -- node[above=1pt,hw label,fill=white,inner sep=1pt]{x=1} (ba2);
  \draw[hw bus] (b2) -- node[above=1pt,hw label,fill=white,inner sep=1pt]{y=1} (bb2);
  % ③ 读旧值
  \node[hw label, anchor=east] at (-0.65,-3.0) {③ 读旧值};
  \node[hw state, minimum width=0.95cm] (a3) at (0,-3.0) {核 A};
  \node[hw control, minimum width=1.0cm] (ba3) at (1.6,-3.0) {x = 1};
  \node[hw memory, minimum width=2.0cm] (m3) at (4.0,-3.0) {x = 0，y = 0};
  \node[hw control, minimum width=1.0cm] (bb3) at (6.4,-3.0) {y = 1};
  \node[hw state, minimum width=0.95cm] (b3) at (8.0,-3.0) {核 B};
  \draw[hw bus] (3.6,-3.35) -- (3.6,-3.8) -- node[pos=0.5,above=1pt,hw label]{r1 = y = 0} (0,-3.8) -- (a3.south);
  \draw[hw bus] (4.4,-3.35) -- (4.4,-4.05) -- node[pos=0.5,below=1pt,hw label]{r2 = x = 0} (8.0,-4.05) -- (b3.south);
  % ④ 屏障后排空
  \node[hw label, anchor=east] at (-0.65,-4.9) {④ 屏障后排空};
  \node[hw state, minimum width=0.95cm] (a4) at (0,-4.9) {核 A};
  \node[hw control, minimum width=1.0cm] (ba4) at (1.6,-4.9) {（空）};
  \node[hw memory, minimum width=2.0cm] (m4) at (4.0,-4.9) {x = 1，y = 1};
  \node[hw control, minimum width=1.0cm] (bb4) at (6.4,-4.9) {（空）};
  \node[hw state, minimum width=0.95cm] (b4) at (8.0,-4.9) {核 B};
  \draw[hw bus] (ba4) -- node[above=1pt,hw label,fill=white,inner sep=1pt]{x=1} (m4);
  \draw[hw bus] (bb4) -- node[above=1pt,hw label,fill=white,inner sep=1pt]{y=1} (m4);
  \node[hw label] at (4.0,-5.75) {mfence 之后：写先全局可见，后续 load 才允许执行};
\end{pdfdiagram}
\diagramnote{store buffer 双核序列的四张快照。x=1、y=1 先滞留在各自核的 store buffer（②），两个 load 直接从内存读回旧值 0（③），于是 r1=0 且 r2=0 成立；mfence 强制 store buffer 排空、写全局可见之后才放行后续 load（④），这个结果即被禁止。四张快照与上表第 1、2、5 步前后的状态一一对应。}

这个结果不违反任何一条单核规则——每个核自己看来程序都按序执行了——但它违反了“所有核看到同一个全局顺序”的直觉。x86 的 \code{mfence} 指令正是为此存在：它强制本核的 store buffer 排空之后，后续的 load 才允许执行，也就是说本核此前的写必须先对其他核可见。

\begin{lstlisting}
core A                    core B
mov qword ptr [x], 1      mov qword ptr [y], 1
mfence                    mfence
mov rax, [y]              mov rbx, [x]
\end{lstlisting}

加入屏障后，r1=0 且 r2=0 被禁止。推理只需直觉：若 A 读到了 y=0，说明 B 的 y=1 在那一刻尚未对 A 可见；而 B 的 \code{mfence} 保证 y=1 全局可见之后 B 才会执行 r2=x，到那时 A 的 x=1 早已因 A 自己的 \code{mfence} 而对所有核可见，所以 B 读 x 必得 1。屏障挡住的正是“写还留在 store buffer 里、读却已经出发”这段窗口；它没有让写变得更快，只是把可见性顺序钉死。

\topic{描述符环把 DMA 变成持续事务流}
高速设备通常不会每次只处理一个缓冲区。网卡、磁盘控制器、USB 控制器常使用描述符环或队列：CPU 准备一批描述符，设备按头尾指针消费，完成后写回状态或 completion 队列，并用中断或轮询通知 CPU。这样可以减少 MMIO 次数，让 CPU 和设备并行工作。

\begin{lstlisting}
receive ring:
  CPU owns descriptor N
  CPU fills buffer address and length
  CPU marks descriptor ready for device
  CPU updates tail pointer or rings doorbell
  device writes packet into buffer
  device writes completion status
  CPU regains ownership and processes buffer
\end{lstlisting}

\begin{longtable}{L{0.18\linewidth}L{0.38\linewidth}L{0.34\linewidth}}
\toprule
环队列字段 & 硬件含义 & 关键问题 \\
\midrule
描述符地址 & 设备要读写的主存位置 & 地址是否经过 IOMMU 或物理地址转换 \\
长度 & 本次 DMA 允许访问的字节数 & 设备不能越界写出缓冲区 \\
所有权位 & 当前由 CPU 还是设备拥有 & 所有权不重叠，状态切换有屏障 \\
头尾指针 & CPU 和设备约定的生产/消费位置 & 指针更新顺序和环绕规则清楚 \\
completion & 设备报告完成、长度、错误码 & 中断到达不等于所有 cache 可见性自动正确 \\
\bottomrule
\end{longtable}

这个模型也说明驱动程序为什么必须理解硬件。驱动不是“调用设备 API”的普通软件，它要写 MMIO 寄存器、分配对齐缓冲区、建立 DMA 映射、维护 cache 可见性、处理中断、回收描述符和处理错误。每一步都对应本章的硬件对象：地址、权限、字节数、所有权、完成和错误。

\topic{描述符环不动、指针动}
描述符环的关键是“环不动、指针动”：描述符数组在内存中的位置固定，CPU 只推进 SW 指针，设备只推进 HW 指针，走到末尾就回绕；两个指针的差值就是环中未完成事务的数量。doorbell 不携带数据，只是告诉设备“指针又动了”，设备即使错过一次 doorbell，靠指针差也能补读。反过来，某个描述符是否完成，不能凭 doorbell 或中断猜测，必须以设备写回描述符状态位（或 completion 队列项）为准；CPU 回收缓冲区前，还要确认这次写回对自己可见。

\begin{pdfdiagram}
  % 两行四格加回绕箭头，代替真圆环，保证正交走线
  \node[hw state, minimum width=2.0cm] (d0) at (1.3,2.6) {描述符 0};
  \node[hw state, minimum width=2.0cm] (d1) at (3.7,2.6) {描述符 1};
  \node[hw state, minimum width=2.0cm] (d2) at (6.1,2.6) {描述符 2};
  \node[hw state, minimum width=2.0cm] (d3) at (8.5,2.6) {描述符 3};
  \node[hw state, minimum width=2.0cm] (d7) at (1.3,1.2) {描述符 7};
  \node[hw state, minimum width=2.0cm] (d6) at (3.7,1.2) {描述符 6};
  \node[hw state, minimum width=2.0cm] (d5) at (6.1,1.2) {描述符 5};
  \node[hw state, minimum width=2.0cm] (d4) at (8.5,1.2) {描述符 4};
  \draw[hw bus, ->] (d0.east) -- (d1.west);
  \draw[hw bus, ->] (d1.east) -- (d2.west);
  \draw[hw bus, ->] (d2.east) -- (d3.west);
  \draw[hw bus, ->] (d3.south) -- (d4.north);
  \draw[hw bus, ->] (d4.west) -- (d5.east);
  \draw[hw bus, ->] (d5.west) -- (d6.east);
  \draw[hw bus, ->] (d6.west) -- (d7.east);
  \draw[hw bus, ->] (d7.north) -- (d0.south) node[midway, fill=white, inner sep=1pt, hw label] {回绕};
  % SW 指针（CPU 写到哪里）与 HW 指针（设备读到哪里）
  \draw[hw return, ->] (8.5,3.6) -- (d3.north);
  \node[hw label, above] at (8.5,3.6) {SW 指针：CPU 下一个写};
  \draw[hw return, ->] (3.7,3.6) -- (d1.north);
  \node[hw label, above] at (3.7,3.6) {HW 指针：设备正在读};
  % doorbell 写动作
  \node[hw control, minimum width=1.7cm] (bell) at (10.65,1.2) {doorbell};
  \draw[hw return, ->] (d3.east) -- (11.05,2.6) -- (11.05,1.51);
  \node[hw label, anchor=east] at (11.0,2.05) {写 doorbell};
\end{pdfdiagram}
\diagramnote{描述符环。设备按 HW 指针顺序消费，CPU 按 SW 指针顺序生产，描述符 7 之后回到 0；CPU 更新 SW 后写一次 doorbell 提示设备。完成与否以设备写回描述符状态位为准，doorbell 只是“有新描述符”的提示。}

\topic{驱动 bring-up 应分层验证}
调试设备时，最差的方法是直接运行完整驱动并等待它“能不能用”。更可靠的顺序是先验证配置空间或设备 ID，再验证 MMIO 映射，再验证单个寄存器读写，再验证中断，再验证一个最小 DMA 描述符，最后才跑持续队列。每一层都要有可观察的结果。

\begin{longtable}{L{0.18\linewidth}L{0.42\linewidth}L{0.3\linewidth}}
\toprule
bring-up 阶段 & 操作 & 预期结果 \\
\midrule
枚举/识别 & 读取设备 ID、版本或能力寄存器 & 地址映射正确，读值稳定可复现 \\
复位 & 写复位位并轮询 ready/busy & 状态位按规格变化，不永久 busy \\
MMIO 回读 & 写可回读配置寄存器再读回 & 字节序、位域和写掩码正确 \\
中断 & 人工触发或配置定时事件 & 中断线/MSI 到达，清除顺序正确 \\
单次 DMA & 一个描述符、一个缓冲区、一个 completion & 地址、长度、cache 可见性和错误码正确 \\
持续队列 & 多描述符环和批量 completion & 所有权、环绕、背压和中断聚合正确 \\
\bottomrule
\end{longtable}

这张表也适用于教学板。即使没有 PCIe，只用一个 GPIO、定时器或简化 DMA 控制器，也应该按同样顺序验证：先确认地址译码，再确认寄存器语义，再确认事件通知，再确认数据搬运。这样本书的硬件 trace 就能从最小 CPU 直接延伸到现代设备驱动。

\topic{案例：一个接收包怎样进入内存}
以网卡接收一个数据包为例。CPU 先分配若干缓冲区，把每个缓冲区的地址和长度写入接收描述符环；随后保证描述符对设备可见，再写 MMIO 门铃或尾指针告诉设备可以接收。设备收到外部数据后，作为 DMA 发起者把数据写入某个缓冲区，再写 completion 或描述符状态，最后触发中断。CPU 处理中断，确认 completion，按一致性规则让缓冲区数据对 CPU 可见，然后把缓冲区交给上层软件。

\begin{longtable}{L{0.14\linewidth}L{0.42\linewidth}L{0.34\linewidth}}
\toprule
阶段 & 硬件事务 & 观察要点 \\
\midrule
准备缓冲区 & CPU 在 RAM 中分配对齐缓冲区 & 物理地址或 I/O 虚拟地址有效，长度足够 \\
写描述符 & CPU STORE 地址、长度、所有权位 & 描述符 cache line 已按平台规则可见 \\
通知设备 & CPU 写 MMIO tail/doorbell & MMIO 不被普通 cache 合并或延迟越界 \\
设备 DMA 写 & 设备发起内存写事务 & IOMMU 权限、字节数、completion 顺序正确 \\
投递完成 & 设备写状态并触发中断 & completion 可读，中断不会早于数据可见 \\
CPU 回收 & 中断处理读取状态和缓冲区 & cache invalidate 或 coherent 规则完成 \\
\bottomrule
\end{longtable}

\begin{lstlisting}
receive path:
  CPU: fill rx_desc[N] with buffer address and length
  CPU: barrier, then MMIO[RX_TAIL] = N
  DEV: reads rx_desc[N]
  DEV: writes packet bytes into buffer
  DEV: writes completion status
  DEV: raises interrupt
  CPU: acknowledges interrupt and reads completion
  CPU: synchronizes buffer visibility
  CPU: consumes packet
\end{lstlisting}

这条路径里至少有三次“完成”。第一，CPU 写 MMIO 门铃完成，只表示设备接口看到了通知；第二，设备 DMA 写缓冲区完成，表示数据已经到达内存系统的某个可见点；第三，中断处理完成，表示 CPU 已经确认状态并回收所有权。把这三者混在一起，是驱动和硬件协同时最常见的错误来源。

\topic{案例：发送路径为什么要区分数据和描述符}
发送数据包时方向相反。CPU 先把待发送数据写入缓冲区，再写发送描述符，最后通知设备。设备读描述符，再 DMA 读缓冲区，把数据发出。这里最关键的可见性顺序是：缓冲区内容必须先对设备可见，描述符必须描述正确地址和长度，门铃必须最后到达。否则设备可能发出旧数据、半写数据或错误长度的数据。

\begin{longtable}{L{0.18\linewidth}L{0.4\linewidth}L{0.32\linewidth}}
\toprule
发送对象 & 可见性要求 & 失败表现 \\
\midrule
数据缓冲区 & CPU 写入的数据先对设备可见 & 设备发送旧包或未初始化字节 \\
描述符 & 地址、长度、标志位完整可见 & 设备读错缓冲区或长度 \\
门铃寄存器 & 在数据和描述符之后对设备可见 & 设备提前启动，读到旧状态 \\
completion & 设备确认已不再读取缓冲区 & CPU 过早重用缓冲区导致数据损坏 \\
\bottomrule
\end{longtable}

发送 completion 到达前，CPU 不应重用缓冲区。即使 MMIO 写门铃已经返回，设备可能还没有读描述符；即使设备已经读描述符，可能还没有读完数据；即使数据已经发出，completion 可能还在队列里等待 CPU 处理。因此所有权位和 completion 队列不是形式主义，而是让 CPU 和设备不同时修改同一块内存的硬件约定。

\topic{案例：块设备读请求的队列化}
磁盘、NVMe、SD 控制器或简化块设备也可以用同一模型理解。CPU 准备命令描述符：逻辑块地址、目标缓冲区、长度、读写方向；设备读取命令后，在内部介质或控制器中完成实际读写，再通过 DMA 把数据写到主存，最后投递 completion。区别只在于设备内部工作可能更慢、队列更深、错误码更复杂。

\begin{longtable}{L{0.18\linewidth}L{0.38\linewidth}L{0.34\linewidth}}
\toprule
块请求字段 & 含义 & 硬件问题 \\
\midrule
LBA 或块号 & 设备内部介质位置 & 越界、坏块、介质错误如何报告 \\
缓冲区地址 & DMA 写入或读取的主存位置 & 对齐、IOMMU、cache 可见性 \\
长度 & 传输字节数或块数 & 不能越过缓冲区和描述符许可范围 \\
方向 & 读设备到内存，或写内存到设备 & DMA 方向决定 cache 维护方式 \\
命令 ID & completion 对应哪个请求 & 队列乱序完成时仍能匹配请求 \\
状态码 & 成功、超时、校验错误、权限错误 & 驱动恢复或上报错误 \\
\bottomrule
\end{longtable}

这个案例能把“设备很慢”讲得更具体。CPU 发出块读请求后，并不会逐字等待数据从设备口读出。它可以切换去做别的工作；设备和 DMA 在后台推进；中断或轮询 completion 时，CPU 再取回结果。硬件性能来自异步队列和批量搬运，而不是让单条 LOAD 直接等磁盘或网络。

\topic{案例：一个未映射访问如何成为可诊断错误}
再看失败路径。CPU 执行 \code{MOV R1,[0xDEAD0000]}，地址译码没有任何目标响应。没有错误约定时，CPU 可能读到悬空总线值，或永久等待 ready。可靠平台应规定超时、总线错误或异常入口。若有 MMU，还可能在页表阶段就发现该地址未映射，不会走到外部总线。

\begin{longtable}{L{0.2\linewidth}L{0.38\linewidth}L{0.32\linewidth}}
\toprule
失败位置 & 可观察现象 & 正确处理 \\
\midrule
页表不存在 & 地址转换阶段触发页错误 & 保存故障地址和错误码，进入页错误入口 \\
权限不允许 & 转换存在但读写/特权不满足 & 触发保护异常，不发起外部访问 \\
地址译码无目标 & 总线事务无人响应 & 超时或 bus error，而不是无限等待 \\
目标返回错误 & 设备或内存控制器报告失败 & 错误状态进入异常或驱动恢复路径 \\
ECC 校验失败 & 数据返回但校验错误 & 更正、重读、上报或机器检查 \\
\bottomrule
\end{longtable}

这张表对硬件调试尤其有用。现象“程序卡住”可能来自页表、总线译码、设备 ready、时钟、复位或错误入口任何一层。把失败位置写成层级，就能决定先看页表项、地址线、片选、ready、错误状态，还是异常向量。

\begin{classicnote}
读任何硬件平台，都可以用同一张事务表约束理解：发起者、目标、地址、方向、数据宽度、握手、完成、错误、副作用和可观察状态。早期 8086 总线、微控制器片内外设、PCIe 设备和现代 SoC 互连的规模不同，但这十项问题不变。
\end{classicnote}

\topic{专题：DRAM 控制器为什么要排队和调度}
SRAM 可以近似看成“地址稳定后读出数据”的存储器，DRAM 则必须先打开行、读写列、预充电、刷新。主存控制器面对的不是单个 LOAD，而是一串来自 CPU、DMA、显示控制器和其他主设备的请求。它要在正确性约束内排序请求，尽量利用行命中和 bank 并行，同时保证刷新不丢失数据。

\begin{longtable}{L{0.18\linewidth}L{0.42\linewidth}L{0.3\linewidth}}
\toprule
DRAM 概念 & 含义 & 对系统的影响 \\
\midrule
row activate & 把某一行搬到行缓冲中 & 同一行后续列访问较快 \\
row conflict & 下一个请求落在不同的行 & 需要预充电再打开新行，延迟增加 \\
bank & 多组可相对独立工作的阵列 & 交错访问可隐藏部分等待 \\
refresh & 周期性恢复电容电荷 & 刷新窗口内部分请求会被延后 \\
ECC & 用冗余位检测或纠正错误 & 增加位宽和延迟，但提高可靠性 \\
\bottomrule
\end{longtable}

因此，主存延迟不是一个固定常数。一次 cache miss 到达内存控制器后，可能遇到行命中、行冲突、刷新、写队列排空或高优先级 DMA。系统软件看到的是“同样的 LOAD 有时快有时慢”，硬件工程师看到的是请求队列、调度策略和时序参数。教学时不必展开 DDR 全部命令，但必须让读者知道：主存不是无限宽、无限快、无排队的数组。

\topic{专题：cache 设计参数怎样改变命中率和时延}
cache 的基本问题是把大而慢的主存切成较小的行，并保存最近使用的数据。地址会被拆成 tag、index、offset：offset 选择 cache 行内字节，index 选择组，tag 判断这一组里哪一路是目标行。不同组织改变冲突概率、比较器数量、替换策略和命中时延。

\begin{longtable}{L{0.18\linewidth}L{0.36\linewidth}L{0.36\linewidth}}
\toprule
组织 & 优点 & 代价或风险 \\
\midrule
直接映射 & 查找简单、速度快、面积小 & 两个常用地址若映到同一行会反复冲突 \\
组相联 & 多路候选降低冲突 miss & 多个 tag 比较器和路选择 mux 增加延迟 \\
全相联 & 任意行可放任意位置，冲突最少 & 硬件代价高，通常用于 TLB 或小 cache \\
写直达 & 内存总是较新，恢复简单 & 写流量大，常需写缓冲 \\
写回 & 多次写合并到 cache 行，带宽效率高 & dirty 行替换时必须写回，失电/一致性更复杂 \\
\bottomrule
\end{longtable}

评价 cache 不能只看“第二次访问变快”。还要看替换、脏行写回、字节写使能、未对齐访问、MMIO 不可 cache、DMA 可见性和异常路径。例如设备把数据 DMA 到主存后，CPU 若继续读旧 cache 行，就会看到过期数据；CPU 写好发送缓冲区后，若脏行没有写回，设备也可能读到旧内容。这就是为什么驱动要有 cache clean/invalidate 或一致性互连。

\topic{专题：MMIO 寄存器位语义要像数据手册一样写}
MMIO 寄存器不是普通变量。每一位都可能有副作用：读清、写 1 清、只读、保留位必须写 0、写后自动启动、硬件置位软件清除。若教材只写“把某地址当寄存器读写”，读者会在真实芯片上犯危险错误。

\begin{longtable}{L{0.18\linewidth}L{0.36\linewidth}L{0.36\linewidth}}
\toprule
位语义 & 例子 & 软件访问要求 \\
\midrule
RO & 输入电平、设备 ID & 写入无效或保留，不应读改写破坏其他位 \\
WO & FIFO 写口、启动命令 & 读值可能无意义，不能依赖回读 \\
RW & 配置模式、使能位 & 修改前要明确复位默认值 \\
W1C & 中断状态位 & 写 1 清除对应事件，写 0 保持 \\
RC & 读状态自动清除 & 调试打印也可能改变设备状态 \\
reserved & 未来扩展或工厂测试位 & 按手册写固定值，通常为 0 \\
\bottomrule
\end{longtable}

一个严谨的 GPIO 寄存器描述至少要写清：方向寄存器决定引脚驱动还是输入；输出寄存器只影响输出模式；输入寄存器来自同步后的外部电平；中断状态由边沿检测置位；清除中断采用 W1C；保留位读值未定义且写 0。这样驱动代码才能避免“读状态时把事件清掉”“读改写时误清中断”“把保留位写成 1 后芯片行为改变”等问题。

\topic{专题：总线错误要能定位到层级}
当一次访问失败时，工程上最重要的问题不是“错了”，而是错在地址转换、权限、互连、目标设备还是数据校验。若每层都只把错误折叠成一个笼统异常，调试会非常困难。更好的设计是保留故障地址、访问方向、主设备 ID、错误来源和目标返回码。

\begin{longtable}{L{0.2\linewidth}L{0.42\linewidth}L{0.28\linewidth}}
\toprule
记录字段 & 作用 & 调试价值 \\
\midrule
fault address & 失败访问的地址 & 区分空指针、越界和 MMIO 映射错误 \\
access type & 读、写、取指、原子操作 & 判断权限和副作用风险 \\
master ID & CPU、DMA、GPU 或调试器 & 找到真正发起者 \\
decode result & 命中哪个目标或未命中 & 定位地址地图错误 \\
response code & OKAY、SLVERR、DECERR、超时等 & 区分目标报错和无人响应 \\
\bottomrule
\end{longtable}

这也是为什么后续平台 bring-up 日志要记录总线层的事务信息。单看 C 代码中的一行 \code{*ptr} 无法知道失败来自页表、cache、IOMMU、互连还是设备时钟未开；完整错误记录能把问题缩小到可测量的硬件边界。

\topic{专题：内存层次要同时看容量、时延、带宽和所有权}
计算机硬件不是只有 CPU。CPU 的每一次有用工作，几乎都要和某一级存储或 I/O 交换信息。寄存器最快但数量极少；L1/L2/L3 cache 容量逐级增大、时延逐级上升；DRAM 容量大但需要控制器、队列和刷新；SSD/硬盘容量更大，却只能通过控制器、队列、DMA 和中断间接参与执行。把这些对象都称为“存储”会掩盖关键差异：有些层按字节随机访问，有些层按 cache line 交换，有些层按页或块传输，有些层只能通过命令队列返回 completion。

\begin{pdfdiagram}[x=1cm,y=1cm]
  \node[hw state, minimum width=1.55cm] (reg) at (0,0) {寄存器};
  \node[hw compute, minimum width=1.55cm] (l1) at (1.8,0) {L1};
  \node[hw compute, minimum width=1.55cm] (l2) at (3.6,0) {L2/L3};
  \node[hw memory, minimum width=1.75cm] (dram) at (5.65,0) {DRAM};
  \node[hw control, minimum width=1.75cm] (ssd) at (7.9,0) {SSD/硬盘};
  \node[hw lane, minimum width=8.8cm, text width=8.4cm] at (3.95,-1.35) {越靠左，访问粒度越小、时延越低、容量越小；越靠右，容量越大，但访问要经过控制器、队列、DMA、错误码和完成通知。};
  \draw[hw bus] (reg) -- node[above,hw label]{操作数} (l1);
  \draw[hw bus] (l1) -- node[above,hw label]{cache line} (l2);
  \draw[hw bus] (l2) -- node[above,hw label]{回填/写回} (dram);
  \draw[hw bus] (dram) -- node[above,hw label]{DMA 缓冲区} (ssd);
\end{pdfdiagram}
\diagramnote{内存层次不是一条“更大的数组”，而是一组访问粒度和完成语义不同的层。}

\begin{longtable}{L{0.16\linewidth}L{0.22\linewidth}L{0.28\linewidth}L{0.24\linewidth}}
\toprule
层次 & 常见粒度 & 谁直接访问 & 主要硬件问题 \\
\midrule
寄存器 & word 或向量 lane & 执行单元 & 端口数、旁路、写回冲突 \\
L1 cache & cache line & 当前核心 load/store 单元 & 命中时延、虚实地址、别名和写策略 \\
L2/L3 cache & cache line & 一个核心或多个核心 & 容量、共享、替换、一致性 \\
DRAM & burst/cache line & 内存控制器 & bank、行命中、刷新、ECC、调度 \\
SSD/硬盘 & 块、页或扇区 & 设备控制器和 DMA & 命令队列、坏块、磨损、错误恢复 \\
网络/显卡等设备 & 描述符、包、帧、buffer & DMA 引擎和设备内部队列 & 所有权、completion、IOMMU、cache 可见性 \\
\bottomrule
\end{longtable}

从程序角度看，\code{load x} 只是取一个变量；从硬件角度看，它可能命中寄存器旁路、命中 L1、命中 L2、命中共享 L3、访问 DRAM、触发页表遍历、发生缺页，甚至等待块设备把页面读回主存。CSAPP 强调的 locality 就是让常见访问停留在左侧短路径。硬件书不能只说“内存很快/很慢”，而要说明每一级的访问粒度、所有者、完成条件和失败路径。

\topic{专题：内存别名与同源不同名的 cache 问题}
上一专题的层次表把“别名”列为 L1 cache 的关键问题之一，这里把它展开。TLB 的存在意味着同一个物理页可以被多个虚拟页映射：共享库、进程间共享内存、同一页在用户态和内核态各有别名，都是日常情形。这带来一个只在“用虚拟地址索引 cache”时才出现的问题：两个不同的虚拟地址指向同一物理行，若 cache 按虚拟地址的低位选槽，而两个别名的虚拟页号不同，index 就可能不同，同一物理行会被存进两个不同的槽。之后程序经别名一写入，经别名二读到的却是另一个槽里的旧副本——同一数据在 cache 里有了两份不一致的拷贝。

看一个小例子。页大小 4 KiB，cache 为 32 KiB 直接映射、64 B 行：offset 为 6 bit，行数 $=32\ \text{KiB}\div 64\ \text{B}=512$，index 取地址的 bit 14..6 共 9 bit。页内偏移只有 12 bit（bit 11..0），因此 index 的高 3 位（bit 12..14）来自虚拟页号——别名地址的差异正好在这几位里。设虚拟地址 \code{0x1040} 和 \code{0x2040} 都映射到物理地址 \code{0x35040}：前者的 index 是 65，后者的 index 是 129，同一物理行会被分别存进槽 65 和槽 129。经 \code{0x1040} 写入后，槽 65 是新值，槽 129 仍是旧值，两次读到同一物理地址却得到不同数据。

两种硬件解法基于同一条思路：让同一物理行的所有别名映到同一个槽，再由物理 tag 确认身份。第一种是物理索引（PIPT）：index 直接取自物理地址，\code{0x1040} 和 \code{0x2040} 经 TLB 翻译后都是 \code{0x35040}，index 都是 321，必然同一个槽，tag 也比较物理地址，别名问题根本不会出现；代价是要先等 TLB 翻译完成才能开始查 cache，命中路径变长。第二种是限制 index 位数（VIPT 的常见约定）：把 cache 设计成 index 加 offset 不超过页内偏移的 12 bit，例如 4 KiB 直接映射，或 32 KiB、8 路组相联——组数 $=32\ \text{KiB}\div 64\ \text{B}\div 8=64$，index 只需 6 bit，index 加 offset 正好 12 bit。此时 index 全部来自页内偏移，而任何别名的页内偏移都与物理地址相同：上例中两个别名与物理地址的 bit 11..6 都等于 1，必然映到同一组。查找与 TLB 翻译并行进行，拿到物理页号后再与组内各路的物理 tag 比较，第二次访问命中同一个行，两份副本无从产生。现代处理器的 L1 大多采用这种“虚拟索引、物理 tag、index 不越出页内偏移”的安排，既保留并行查找的速度，又让别名退化为普通命中。把两种安排画成对照图，差别就看得很清楚。

\begin{pdfdiagram}
  % 左：未限制 index 位数（32 KiB 直接映射，index 取 bit 14..6）
  \node[hw label] at (1.9,2.75) {未限制：按虚拟 index 映到两个槽};
  \node[hw state, minimum width=1.6cm] (va1) at (-0.3,0.9) {VA \code{0x1040}};
  \node[hw state, minimum width=1.6cm] (va2) at (-0.3,-0.9) {VA \code{0x2040}};
  \node[hw compute, minimum width=1.0cm] (tlb) at (1.8,0) {TLB};
  \node[hw state, minimum width=1.9cm] (pa) at (3.95,0) {PA \code{0x35040}};
  \node[hw memory, minimum width=2.0cm] (s65) at (3.95,1.9) {槽 65：新值};
  \node[hw memory, minimum width=2.0cm] (s129) at (3.95,-1.9) {槽 129：旧值};
  \draw[hw bus] (va1.north) -- (-0.3,1.9) -- node[pos=0.7,above=1pt,hw label,fill=white,inner sep=1pt]{index 65} (s65.west);
  \draw[hw bus] (va2.south) -- (-0.3,-1.9) -- node[pos=0.7,above=1pt,hw label,fill=white,inner sep=1pt]{index 129} (s129.west);
  \draw[hw bus] (va1.east) -| (tlb.north);
  \draw[hw bus] (va2.east) -| (tlb.south);
  \draw[hw bus] (tlb) -- node[above=1pt,hw label,fill=white,inner sep=1pt]{同一 PA} (pa);
  \draw[hw return] (pa.north) -- node[right=1pt,hw label]{tag} (s65.south);
  \draw[hw return] (pa.south) -- node[right=1pt,hw label]{tag} (s129.north);
  \node[hw label] at (1.9,-2.75) {同一物理行出现两份副本};
  % 右：限制 index 位数（index+offset 不越出页内偏移 12 bit）
  \node[hw label] at (7.3,2.75) {限制 index 位数：别名同组};
  \node[hw state, minimum width=1.9cm] (var) at (6.2,0.9) {VA \code{0x1040}\\VA \code{0x2040}};
  \node[hw compute, minimum width=1.0cm] (tlb2) at (8.3,0.9) {TLB};
  \node[hw memory, minimum width=2.0cm] (grp) at (8.3,-0.9) {组 1：唯一副本};
  \draw[hw bus] (var) -- (tlb2);
  \draw[hw bus] (var.south) -- node[pos=0.4,right=1pt,hw label]{index = bit 11..6} (6.2,-0.9) -- (grp.west);
  \draw[hw return] (tlb2.south) -- node[right=1pt,hw label]{PA \code{0x35040}\\物理 tag} (grp.north);
  \node[hw label] at (7.3,-2.75) {index 全来自页内偏移};
\end{pdfdiagram}
\diagramnote{VIPT 别名问题与限制 index 位数的对照。左：32 KiB 直接映射 cache 的 index 取 bit 14..6，\code{0x1040} 与 \code{0x2040} 的 index 分别是 65 与 129，同一物理行 \code{0x35040} 被存进两个槽，经 \code{0x1040} 写入后槽 65 是新值、槽 129 仍是旧值。右：限制 index 加 offset 不越出页内偏移 12 bit（如 32 KiB、8 路组相联，index 为 bit 11..6），两个别名的 index 都来自页内偏移，必然映到同一组，再由物理 tag 确认身份，双副本无从产生。}

\topic{专题：cache 与 TLB 是两套不同的近路}
cache 缓存数据或指令内容，TLB 缓存虚拟页到物理页的地址转换。二者经常同时出现在一次 LOAD 路径上，但回答的问题不同：TLB 问“这个虚拟地址能否翻译成哪个物理地址”；cache 问“这个物理位置的内容是否已经在近处”。把二者混在一起，会导致读者无法理解为什么 TLB miss 不等于 cache miss，页错误也不等于总线错误。

\begin{pdfdiagram}
  \node[hw state, minimum width=1.85cm] (va) at (0,0) {虚拟地址};
  \node[hw compute, minimum width=1.85cm] (tlb) at (2.45,0) {TLB};
  \node[hw state, minimum width=1.85cm] (pa) at (4.9,0) {物理地址};
  \node[hw compute, minimum width=1.85cm] (cache) at (7.35,0) {cache};
  \node[hw memory, minimum width=1.65cm] (data) at (9.65,0) {数据};
  \node[hw control, minimum width=2.0cm] (ptw) at (2.45,-1.35) {页表遍历};
  \node[hw memory, minimum width=2.0cm] (mem) at (7.35,-1.35) {DRAM 回填};
  \draw[hw bus] (va) -- (tlb);
  \draw[hw bus] (tlb) -- node[above,hw label]{命中} (pa);
  \draw[hw bus] (pa) -- (cache);
  \draw[hw bus] (cache) -- node[above,hw label]{命中} (data);
  \draw[hw danger] (tlb) -- node[right,hw label]{miss} (ptw);
  \draw[hw danger] (cache) -- node[right,hw label]{miss} (mem);
\end{pdfdiagram}
\diagramnote{TLB miss 发生在地址转换阶段，cache miss 发生在数据回填阶段。二者都可能让 LOAD 变慢，但失败位置不同。}

\begin{longtable}{L{0.2\linewidth}L{0.34\linewidth}L{0.36\linewidth}}
\toprule
事件 & 真实含义 & 后续动作 \\
\midrule
TLB hit + cache hit & 地址转换和数据都在近处 & LOAD 快速返回 \\
TLB miss + 页表命中 & 转换不在 TLB，但页表有效 & 硬件或软件页表遍历后填 TLB \\
页不存在 & 页表没有有效映射 & 进入页错误异常，可能由 OS 分配或从磁盘读页 \\
cache miss & 地址有效，但数据不在当前 cache & 从下一级 cache 或 DRAM 回填 cache line \\
权限错误 & 转换存在但访问类型不允许 & 报保护异常，不能提交数据 \\
\bottomrule
\end{longtable}

\topic{专题：SSD 和硬盘不是慢内存，而是块设备}
硬盘和 SSD 都不能像 SRAM 那样把地址线接上就读一个字节。CPU 要通过控制器提交命令：读哪个逻辑块、读多长、放到哪个主存缓冲区、完成后写哪个 completion。硬盘的内部约束来自机械寻道和旋转；SSD 的内部约束来自 NAND flash 页读写、块擦除、FTL 映射、磨损均衡、垃圾回收和掉电保护。它们对 CPU 暴露的共同形态，是队列化块设备。

\begin{longtable}{L{0.18\linewidth}L{0.34\linewidth}L{0.38\linewidth}}
\toprule
对象 & 内部结构重点 & 对系统可见的边界 \\
\midrule
机械硬盘 & 盘片、磁头、磁道、扇区、缓存 & 随机访问受寻道和旋转影响，顺序访问更友好 \\
SATA SSD & NAND flash、FTL、通道、擦除块 & 写放大、垃圾回收、坏块管理、掉电保护 \\
NVMe SSD & PCIe、多队列、doorbell、completion queue & 并发队列深度高，CPU 通过 MMIO 和 DMA 协作 \\
SD/eMMC & 控制器隐藏 flash 细节，命令集较简单 & 嵌入式启动和存储，延迟波动明显 \\
\bottomrule
\end{longtable}

理解 SSD 时要避免两个误解。第一，SSD 没有机械寻道，不等于每次访问都一样快；内部垃圾回收、擦除、映射表缓存和热控都会让尾延迟上升。第二，逻辑块地址连续，不等于物理 NAND 页连续；FTL 会把逻辑写重映射到新位置，再在后台回收旧块。因此块设备的正确抽象不是“巨大数组”，而是“命令队列 + DMA 缓冲区 + completion + 错误恢复”。

\begin{pdfdiagram}
  \node[hw state, minimum width=1.55cm] (cpu) at (0,0) {CPU};
  \node[hw control, minimum width=2.05cm] (sq) at (2.55,0.72) {提交队列};
  \node[hw control, minimum width=1.65cm] (bell) at (2.55,-0.72) {门铃};
  \node[hw control, minimum width=2.0cm] (ctrl) at (5.25,0) {NVMe 控制器};
  \node[hw memory, minimum width=1.85cm] (nand) at (7.8,0.72) {NAND/FTL};
  \node[hw state, minimum width=2.05cm] (cq) at (7.8,-0.72) {完成队列};
  \node[hw label] at (1.25,1.15) {1 写命令};
  \node[hw label] at (1.25,-1.12) {2 敲门铃};
  \node[hw label] at (6.55,-1.18) {3 投递完成};
  \draw[hw bus] (cpu) -- (sq);
  \draw[hw bus] (cpu) -- (bell);
  \draw[hw bus] (sq) -- (ctrl);
  \draw[hw bus] (bell) -- (ctrl);
  \draw[hw bus] (ctrl) -- (nand);
  \draw[hw return] (ctrl) -- (cq);
  \draw[hw return] (cq.west) |- (cpu.south);
\end{pdfdiagram}
\diagramnote{NVMe 的基本形态：CPU 写队列和门铃，设备 DMA 读写主存，completion 才说明请求完成。}

\topic{专题：读硬件结构时要区分数据路径和控制路径}
CPU、cache、DRAM、SSD、网卡和显卡都可以用同一个提问框架阅读：数据从哪里到哪里，控制由谁发起，完成由谁报告，错误由谁保存。低质量描述常把这些混成“设备把数据给 CPU”。更好的描述应写出两条路径：数据可能经 DMA 在设备和主存之间移动；控制通常由 CPU 写 MMIO、设备写 completion、CPU 读状态或处理中断。

\begin{longtable}{L{0.18\linewidth}L{0.36\linewidth}L{0.36\linewidth}}
\toprule
路径 & 典型方向 & 必须说明 \\
\midrule
控制路径 & CPU 到 MMIO/doorbell，设备到状态/中断 & 命令何时被接收，完成如何通知 \\
数据路径 & 设备 DMA 与主存缓冲区之间 & 缓冲区地址、长度、cache 可见性、权限 \\
错误路径 & 设备或互连到错误码/异常 & 超时、校验、越界、权限和重试策略 \\
恢复路径 & 驱动或固件到 reset/retry/rebuild & 哪些状态保留，哪些队列丢弃或重放 \\
\bottomrule
\end{longtable}

\section*{章末练习}

\begin{longtable}{L{0.22\linewidth}L{0.5\linewidth}L{0.18\linewidth}}
\toprule
任务 & 要求 & 预期结果 \\
\midrule
SRAM 读写周期 & 画出地址、CS、OE、WE、数据线和采样窗口 & 能指出谁驱动数据线，谁采样 \\
地址映射 & 设计 64 KiB ROM/RAM/GPIO/定时器映射 & 任意地址最多一个片选有效 \\
等待状态 & 为慢 ROM 和慢 I/O 设计 READY/WAIT 规则 & 等待期间地址和方向稳定，有超时策略 \\
MMIO 寄存器 & 设计 GPIO 的 \code{DIR/OUT/IN/STAT} 语义 & 外部输入经过同步，状态清除规则明确 \\
中断顺序 & 写出设备 ready 到 CPU 服务程序返回的 trace & 不丢事件，不重复清旧事件 \\
DMA 顺序 & 写出描述符、缓冲区、completion 和 cache 可见性边界 & CPU 和设备所有权不重叠 \\
valid/ready & 用 valid/ready 描述一次读请求和读响应 & 能说明背压和 completion 的作用 \\
cache 行 & 用 tag/index/offset 解释一次 cache 命中和一次 miss & 能说明 cache line、valid、dirty 和替换边界 \\
MMIO 属性 & 说明为什么设备寄存器通常不能按普通 cacheable 内存处理 & 能指出副作用、顺序和可见性要求 \\
未映射访问 & 规定未映射地址的超时、错误或固定返回值 & CPU 不会永久等待未知目标 \\
存储层次图 & 画出寄存器、L1/L2/L3、DRAM、SSD/硬盘的数据粒度和完成语义 & 能区分 cache line、页、块和 completion \\
TLB/cache 区分 & 写出一次 LOAD 中 TLB hit/miss、页错误、cache miss 的不同路径 & 不把地址转换错误和数据回填混为一谈 \\
块设备队列 & 为 NVMe 或简化 SSD 写出 submission、doorbell、DMA、completion 顺序 & 能说明 CPU 写门铃不等于 I/O 完成 \\
DRAM 时序计算 & 给定 tRCD、tCL、时钟周期和突发长度，估算一次随机读时延，并与行命中情形对比 & 三段时延拆分正确，量级为几十 ns \\
UART 帧分析 & 给定 RX 波形和标称波特率，数出起始位、数据位、停止位，还原字节并核对波特率误差 & 帧结构和误差判断都有依据 \\
组相联查找 & 用本章的四次访问序列，分别在直接映射和 2 路组相联下走一遍 & 能指出哪几次由 miss 变 hit 及原因 \\
描述符环顺序 & 写出 CPU 提交两个描述符到设备完成的 SW/HW 指针、doorbell 与完成写回顺序 & doorbell 不当作完成，写回可见性明确 \\
bank 交错拍数 & 4 个 bank 按 cache line 轮转交错，单 bank 一次访问占用 4 拍，bank 间每拍可启动一个新请求；计算连续读 8 条相邻 cache line 的总拍数，并与单 bank 串行对比 & 交错 11 拍、单 bank 32 拍，能写出每行的启动与完成拍 \\
替换策略 & 2 路组相联、LRU，某组初始为空，依次访问 A、B、A、C、A、B；逐步写出命中/失效和组内容，并与同 index 直接映射对比 & 2 路共 3 次 miss，直接映射 6 次全 miss \\
预取演算 & 顺序读 16 条 cache line，每次 miss 延迟 30 拍、每行处理 8 拍；计算无预取与 next-line 预取的总拍数，并求能完全隐藏预取的最小处理节拍 & 608 拍对 488 拍，处理节拍不小于 30 才能完全隐藏 \\
I2C 事务序列 & 写出主机从 7 位地址 \code{0x50} 的从机寄存器 \code{0x10} 读 2 字节的完整序列：START、地址字节、读写位、ACK/NACK、repeated START、STOP & \code{0xA0}/\code{0xA1} 正确，最后一字节 NACK，repeated START 位置正确 \\
仲裁对照 & 用一句话分别说明集中式仲裁器、菊花链、轮询计数的线数、优先级和公平性，并指出片上互连中仲裁的位置 & 三种做法的取舍清楚，仲裁做进交换节点 \\
MESI 走查 & 两核初始均为 I，按“A 读 X、B 读 X、A 写 X、B 读 X”写出每步两核状态和互连消息 & 每步最新值位置唯一，失效与写回/转发正确 \\
\bottomrule
\end{longtable}

\section*{参考解答与要点}

\begin{longtable}{L{0.18\linewidth}L{0.48\linewidth}L{0.24\linewidth}}
\toprule
任务 & 参考解答 & 必须保留的要点 \\
\midrule
SRAM 周期 & 读时地址和片选稳定，OE 有效、WE 无效，由 SRAM 驱动数据；写时 OE 无效，由发起者驱动数据，WE 在地址和数据稳定后有效。 & 输出使能互斥，写窗口满足建立/保持时间。 \\
地址映射 & 例如低半区 ROM，高半区中一块 RAM，最高 4 KiB 为 I/O；用高位地址产生互斥片选，未映射区返回错误或超时。 & 片选互斥和未映射行为必须写入规格。 \\
等待状态 & 慢目标在 ready 之前保持事务未完成，控制器保持地址、片选和方向；超过最大等待则返回错误。 & 等待状态是硬件事务延长。 \\
MMIO & GPIO 输出写入锁存器，输入来自同步后的外部引脚，状态位说明 ready/error/irq，清除规则避免新旧事件混淆。 & MMIO 有副作用，不是普通变量。 \\
中断 & 设备锁存事件并请求中断，中断控制器投递，CPU 进入入口，服务程序读状态、处理数据、按设备规则清事件，再发送 EOI 或恢复中断。 & 清除顺序决定是否丢事件。 \\
DMA & CPU 准备描述符和缓冲区并保证对设备可见，设备搬运数据并投递 completion，CPU 在 completion 后按一致性规则读取缓冲区。 & 描述符可见、数据可见和缓冲区所有权必须分开。 \\
valid/ready & 只有 valid 和 ready 同时为 1 才传输；ready 为 0 时发送方保持 payload；读请求和读响应可隔多个周期。 & 背压不等于错误，completion 才是完成标志。 \\
未映射访问 & 未映射区域可以返回固定值、产生错误或等待到超时，但规则必须固定且可观察；不能让总线悬空成为随机值。 & 错误路径也是硬件规格。 \\
cache 行 & 先算 offset=log2（行大小）、index=log2（行数）、tag=剩余位；命中要求槽 valid 且 tag 相等；miss 时取整行填入并更新 tag/valid，dirty 位决定替换时是否要先写回下一级。 & 三个字段的分工和 valid/dirty 语义不能丢。 \\
MMIO 属性 & 设备寄存器有副作用：写是命令、读可能清状态；若被 cache 缓冲或合并，设备可能永远看不到命令，CPU 也可能一直读到旧状态，所以设备内存要映射为不可缓存、不可合并、顺序保持。 & 副作用、顺序、可见性三条都要答到。 \\
存储层次图 & 寄存器（字，当拍完成）、L1/L2/L3（cache line，命中即返回）、DRAM（行/页，控制器调度与刷新）、SSD/硬盘（块，I/O 请求和 completion）；越往下粒度越大，完成语义越依赖显式通知而不是当拍返回。 & 粒度与完成语义要对应，不能只画容量和速度。 \\
TLB/cache 区分 & TLB miss 发生在地址转换阶段：硬件查页表，缺页则走页错误异常；cache miss 发生在拿到物理地址之后：向下一级回填整行，CPU 只是停顿或乱序等待。页错误是控制流改道，cache miss 只是数据慢。 & 转换错误与数据回填是两条独立路径。 \\
块设备队列 & CPU 把描述符写进提交队列，写 doorbell 通知设备；设备 DMA 取走描述符、搬完数据后写完成队列，再（可选）发中断。doorbell 只表示“有新请求”，不代表 I/O 完成，完成以完成队列项为准。 & submission 与 completion 分离，doorbell 不等于完成。 \\
DRAM 时序 & 一次需要开行的随机读约为 tRCD（开行）加 tCL（列出数）加突发传输，常见 DDR3/DDR4 器件合计约 30--40 ns 量级；行命中可省去 tRCD，行冲突还要先 precharge。 & 时延由行、列、突发三段组成，不是固定常数。 \\
UART 帧 & 先找起始位下降沿，再按位时间中点逐位采样：1 起始位、8 数据位（低位在前）、1 停止位；还原字节后核对实测位时间与标称波特率的偏差，双方每一侧误差一般不超过约 2\%，收发合计留有约 5\% 的理论上限。 & 位边界、采样点和误差都要核对。 \\
组相联 & index 选中整个组，组内两路 tag 并行与地址 tag 比较，任一路相等且 V=1 即 hit；两个热地址各占一路后不再互相驱逐；miss 时按替换位选一路驱逐。 & 并行比较、mux 和替换位缺一不可。 \\
描述符环 & CPU 只推进 SW 指针，设备只推进 HW 指针，末尾回绕；doorbell 只提示有新描述符，不代表完成；完成以设备写回状态位或 completion 项为准，回收前确认写回可见。 & 指针分离、doorbell 非完成、写回可见性。 \\
bank 交错 & 第 1..4 拍依次启动第 1..4 行（各占一个 bank），第 $i$ 行在第 $i$ 拍启动、第 $i+3$ 拍完成；第 5 行在第 5 拍回到 bank 1 时前一次访问刚结束，无需等待。第 8 行在第 11 拍完成，共 11 拍。单 bank 串行需 $8\times4=32$ 拍。 & 加速来自启动流水化，单 bank 本身没有变快。 \\
替换策略 & 2 路 LRU：A miss 得 \{A,--\}；B miss 得 \{A,B\}；A hit（B 成为最久未用）；C miss 驱逐 B 得 \{A,C\}；A hit；B miss 驱逐 C 得 \{A,B\}，共 3 次 miss。直接映射下同 index 只有一个槽，每次访问都驱逐上一行，6 次全 miss。 & LRU 驱逐“最久未用”，不是“最早进入”。 \\
预取演算 & 无预取：每行 miss 30 拍加处理 8 拍，$16\times38=608$ 拍。next-line 预取：第 1 行第 30 拍就绪，处理第 $i$ 行的 8 拍只能隐藏预取的一部分，第 $i$ 行数据在第 $30i$ 拍就绪，稳态每行 30 拍，总 $30\times16+8=488$ 拍。处理节拍 $T\geq30$ 时预取被完全隐藏，总拍数降为 $30+16T$。 & 预取隐藏延迟但不消除首次 miss；处理太慢时仍受 miss 间隔限制。 \\
I2C 事务 & START $\rightarrow$ \code{0xA0}（\code{0x50} 左移一位、读写位 0 表示写）$\rightarrow$ 从机 ACK $\rightarrow$ \code{0x10}（寄存器号）$\rightarrow$ ACK $\rightarrow$ repeated START $\rightarrow$ \code{0xA1}（读写位 1 表示读）$\rightarrow$ ACK $\rightarrow$ 读第 1 字节 $\rightarrow$ 主机 ACK $\rightarrow$ 读第 2 字节 $\rightarrow$ 主机 NACK $\rightarrow$ STOP。 & 先写寄存器地址，repeated START 转读；最后字节 NACK 通知从机结束。 \\
仲裁对照 & 集中式：每主设备一对请求/授权线，仲裁器并行比较，最快但线数为 $2N$ 且有单点；菊花链：授权逐级传递，线数与设备数无关，但优先级由位置固定、传播慢；轮询计数：$\lceil\log_2 N\rceil$ 根扫描线，从上次授权点继续则轮转公平，但平均扫半表最慢。片上互连把仲裁做进每个交换节点的端口。 & 三者在线数、延迟、公平性之间取舍；仲裁没有消失。 \\
MESI 走查 & 步骤 1：A 由 I 变 S，B 保持 I（读入共享副本）；步骤 2：B 由 I 变 S，A 保持 S；步骤 3：A 广播失效，B 由 S 变 I，A 由 S 变 M，内存过期；步骤 4：B 读 miss，A 写回或转发后由 M 降为 S，B 得 S，内存恢复最新。 & 每步最新值唯一：S 期间在内存，M 期间在 A 的 cache。 \\
\bottomrule
\end{longtable}
